\documentclass[
  reprint,
  amsmath,
  amssymb,
  aps,
  pra,
  nofootinbib,
  superscriptaddress,
  floatfix
]{revtex4-2}

\usepackage[T1]{fontenc}
\usepackage{lmodern}
\usepackage{graphicx}
\graphicspath{{../../}{}}
\usepackage{capt-of} % \captionof for non-floating figures (used with widetext)
\usepackage{booktabs}
\usepackage{bm}
\usepackage{mathtools}
\usepackage{amsthm}
\usepackage{xcolor}
\usepackage{hyperref}

\newcommand{\placeholder}[1]{\textcolor{red!70!black}{\ifmmode\left[#1\right]\else\textsf{[#1]}\fi}}
\newcommand{\tentative}[1]{\textcolor{red!70!black}{#1}}
\newcommand{\structlabel}[1]{\noindent\mbox{\tentative{#1.}}~}
\newcommand{\Oset}{\mathcal{O}}
\newcommand{\Gset}{\mathcal{G}}
\newcommand{\Rset}{\mathcal{R}}
\newcommand{\Sset}{\mathcal{S}}
\newcommand{\Hcal}{\mathcal{H}}
\newcommand{\Bcal}{\mathcal{B}}
\newcommand{\Ocal}{\mathcal{O}}
\newcommand{\Ucal}{\mathcal{U}}
\newcommand{\Jcal}{\mathcal{J}}
\newcommand{\Pcal}{\mathcal{P}}
\newcommand{\Dcal}{\mathcal{D}}
\newcommand{\Mcal}{\mathcal{M}}
\newcommand{\Wcal}{\mathcal{W}}
\newcommand{\dd}{\mathrm{d}}
\newcommand{\ii}{\mathrm{i}}
\newcommand{\ee}{\mathrm{e}}
\newcommand{\TR}{\mathrm{TR}}
\newcommand{\ex}{\mathrm{ex}}
\newcommand{\ctrl}{\mathrm{ctrl}}
\newcommand{\stat}{\mathrm{stat}}
\newcommand{\dyn}{\mathrm{dyn}}
\newcommand{\drv}{\mathrm{drv}}
\newcommand{\obs}{\mathrm{obs}}
\newcommand{\stay}{\mathrm{stay}}
\newcommand{\app}{\mathrm{append}}
\newcommand{\ablate}{\mathrm{ablate}}
\newcommand{\Var}{\operatorname{Var}}
\newcommand{\Tr}{\operatorname{Tr}}
\newcommand{\Top}{\operatorname{Top}}
\newcommand{\argminp}{\mathop{\mathrm{arg\,min}}}
\newcommand{\argmaxp}{\mathop{\mathrm{arg\,max}}}
\newcommand{\proj}{\mathrm{proj}}
\newcommand{\diag}{\mathrm{diag}}
\newcommand{\CVaR}{\operatorname{CVaR}}
\newcommand{\softmax}{\operatorname{softmax}}
\newcommand{\hh}{\mathrm{HH}}
\newcommand{\phmax}{n_{\mathrm{ph,max}}}
\newcommand{\twq}{\mathrm{2Q}}
\newcommand{\eps}{\varepsilon}
\newcommand{\ket}[1]{|#1\rangle}
\newcommand{\bra}[1]{\langle #1|}
\newcommand{\braket}[2]{\langle #1|#2\rangle}
\newcommand{\expval}[1]{\langle #1\rangle}
\newcommand{\inlinetablecaption}[2]{%
  \refstepcounter{table}\label{#1}%
  \noindent\textsc{Table~\thetable.} #2\par\vspace{0.6ex}%
}
\newenvironment{schurlemma}[1][]
{\par\medskip\noindent\textbf{Lemma}\if\relax\detokenize{#1}\relax\else\ (#1)\fi.\space\itshape}
{\par\medskip}
\newenvironment{schurassumption}[1][]
{\par\medskip\noindent\textbf{Assumption}\if\relax\detokenize{#1}\relax\else\ (#1)\fi.\space}
{\par\medskip}
\newenvironment{schurtheorem}[1][]
{\par\medskip\noindent\textbf{Theorem}\if\relax\detokenize{#1}\relax\else\ (#1)\fi.\space\itshape}
{\par\medskip}
\newenvironment{schurremark}[1][]
{\par\medskip\noindent\textbf{Remark}\if\relax\detokenize{#1}\relax\else\ (#1)\fi.\space}
{\par\medskip}


\begin{document}

\title{Geometry- and Cost-Aware ADAPT Ansatz Construction for Mixed Fermion--Boson Hamiltonians}

\author{Jake Strobel}
\author{Jacopo Simoni}
\author{Gabrille Riva}
\author{Yuan Ping}
\affiliation{University of Wisconsin--Madison}

\date{\today}


\begin{abstract}
Mixed fermion--boson Hamiltonians combine electronic correlation with truncated bosonic motion, producing compact few-site benchmarks whose informative observables are already difficult for noisy quantum processors because two-qubit depth, measurement burden, and variational refit cost grow rapidly.  We present SNAKE, a geometry- and cost-aware ADAPT procedure for ground-state ansatz construction.  Candidate generators are promoted to candidate--position records and scored through lower-confidence Fubini--Study trust-region gain, hardware and measurement cost penalties, tangent novelty, reduced-window Schur reranking, batching, beam continuation, and rollback-safe generator ablation.  Across Hubbard--Holstein and spin--boson benchmarks, SNAKE is evaluated against fixed ansatz VQE, append-only ADAPT, Qubit-ADAPT-VQE, TETRIS-ADAPT-VQE, and Geo-ADAPT-VQE.  Relative to the lowest-error recovered non-SNAKE competitor in the fermion--boson aggregate, SNAKE lowers mean energy error from \tentative{$0.2147$} to \tentative{$4.89\!\times\!10^{-5}$}, a \tentative{$99.98\%$} reduction, with the compiled-resource columns reporting the corresponding Pareto tradeoff.
\end{abstract}

\maketitle


\section{Introduction}

\structlabel{Problem, Hamiltonian, importance of solving it}
Mixed fermion--boson Hamiltonians encode electronic correlation, lattice motion, and electron--phonon dressing within a single truncated Hilbert space \cite{Hubbard1963,Holstein1959a,Holstein1959b}.  The problem is to prepare low-energy states of these models at a hardware cost compatible with present noisy processors, while retaining enough bosonic expressivity to resolve density, doublon, site-occupation, and phonon-dressing observables.  Classical references remain essential, but their costs set the motivation for a circuit-level Pareto problem: dense exact diagonalization scales as $O(D_{\rm trunc}^3)$ in the truncated Hilbert dimension \cite{Weisse2008}, sparse eigensolvers still scale with Hamiltonian sparsity and the requested low-energy subspace \cite{Schollwock2011}, fermionic Monte Carlo can encounter sign-problem scaling \cite{Troyer2005}, and tensor-network cost can grow rapidly with bond dimension and bosonic local dimension \cite{Schollwock2011}.  Quantum algorithms move this burden rather than removing it.  A fixed ansatz buys a predetermined depth and optimizer dimension \cite{Kandala2017}; append-only ADAPT screens an operator pool at each iteration \cite{Grimsley2019}; and insertion-aware ADAPT expands the raw acquisition surface from $O(|\Gset|)$ tail candidates to $O(\sum_{m\in\Gset}|\Pcal_m|)$ candidate--position records before any metric, refit, or measurement decision is made \cite{Anastasiou2024,Ramoa2025}.  We consider Hamiltonians with the generic decomposition \cite{Dutta2024BosonicChemistry,Peng2025BosonHamiltonians}
\begin{align}
H&=H_f+H_b+H_{fb},
\nonumber\\
H_f&=\sum_{pq}h_{pq}c_p^\dagger c_q
+\frac12\sum_{pqrs}V_{pqrs}c_p^\dagger c_q^\dagger c_s c_r,
\nonumber\\
H_b&=\sum_{\mu}\omega_\mu b_\mu^\dagger b_\mu+H_b^{\rm nl},
\nonumber\\
H_{fb}&=\sum_{\mu}X_\mu F_\mu[c^\dagger,c]
+\sum_{\mu}P_\mu G_\mu[c^\dagger,c],
\label{eq:general_fb_terms_static}
\end{align}
where $X_\mu=b_\mu^\dagger+b_\mu$ and $P_\mu=\ii(b_\mu^\dagger-b_\mu)$.  Hubbard, Hubbard--Holstein, boson-only, spin-boson, molecular electronic, and vibronic instances are treated as limiting cases. Detailed benchmark families can be found in Appendix~\ref{app:hamiltonian_families}.  Specific phonon-boson qubit encodings can be found in Appendix~\ref{app:encodings}, which have been studied in digital and hybrid quantum-simulation settings \cite{Macridin2018,Sawaya2020,Li2023,Denner2023}.  Because hardware noise scales with circuit depth, entangling operations, routing overhead, and measurement burden, practical NISQ optimization revolves around an accuracy--resource Pareto trade-off between minimizing circuit cost and ensuring simulation expressibility.

\structlabel{Current solution strategies}
NISQ ground-state algorithms usually minimize a Rayleigh--Ritz objective on a parameterized circuit \cite{Peruzzo2014,McClean2016,Kandala2017},
\begin{equation}
E(\theta)=
\bra{\psi_{\rm ref}}U^\dagger(\theta)HU(\theta)\ket{\psi_{\rm ref}},
\label{eq:rayleigh_static}
\end{equation}
where \(\ket{\psi_{\mathrm{ref}}}\) is typically a Hartree--Fock determinant \(\ket{\mathrm{HF}}\); for mixed systems with phonon truncation \(n_{\mathrm{cut}}\), we set \(\ket{\psi_{\mathrm{ref}}}=\ket{\mathrm{HF}}\otimes\ket{0_1\,0_2\,\cdots\,0_{n_{\mathrm{cut}}}}\).  Fixed-structure VQE chooses \(U(\theta)\) in advance \cite{Peruzzo2014,McClean2016,Kandala2017}, whereas ADAPT methods iteratively grow the circuit by appending generators from a pool according to local energy gradients \cite{Grimsley2019,Tang2021,Yordanov2021}. In either case, the parameterization update rule is
\begin{equation}
\label{eq:theta_opt}
\theta^\star=\arg\min_\theta E(\theta).
\end{equation}
Recent variants improve circuit compactness by batching compatible generators (TETRIS-ADAPT) \cite{Anastasiou2024}, using operator sequencing with coupled-exchange strategies (CEO-ADAPT) \cite{Ramoa2025}, and promoting novel tangent directional growth (Overlap-ADAPT) \cite{Feniou2023}. Although traditional ADAPT uses first-order gradient information to expand the circuit, second-order information improves convergence \cite{Grimsley2023ADAPT} and allows the warm-start of parameters across iterations to cut optimizer overhead \cite{Ramoa2025HessianRecycling}. Concurrent with this work, Geo-ADAPT-VQE further improved convergence by replacing the \(\theta\)-space Euclidean metric with a state-space Fubini--Study metric within Eq.~(3) operator-selection score, and enabled generator insertion instead of append-only \cite{Sohail2026GeoADAPT}. In a similiar vein, Prune ADAPT tracks plaussibly archaic generators to ablate based off amplitude and iteration history \cite{Vaquero2025PrunedADAPT}.

\structlabel{Problem with current solutions}
The technical difficulty is the live ranking of a candidate--position record $r=(m,p)$, where inserting generator $A_m$ at position $p$ induces the dressed tangent $U_{>p}A_mU_{>p}^{\dagger}\ket{\psi}$.  A raw ADAPT score is a unit-parameter directional signal; it does not normalize by Fubini--Study state displacement, test overlap with the current tangent span, check whether a local refit window can already recover the direction, or price the candidate's compiled and measurement footprint.  Two records with comparable raw gradients can still differ greatly in two-qubit count, routing-sensitive depth, Pauli support, grouped-measurement burden, and shot requirements.  Insertion position also matters because noncommuting generators alter their state-space action when moved through the ordered ansatz.  Meanwhile, full refits become more expensive with the active parameter count, and stale generators continue to contribute depth unless deletion is confirmed by a recoverability test.  Broad variational landscapes and gradient-driven adaptive searches can also exhibit trainability and selection pathologies \cite{McClean2018BarrenPlateaus,Stadelmann2025GradientTroughs}.  The selector therefore has to optimize post-refit state-space gain against compiled, measurement, and optimization cost, rather than maximizing raw gradient magnitude alone.

\structlabel{Solution strategy}
We introduce SNAKE, a geometry- and cost-aware ADAPT procedure that ranks candidate--position records rather than bare generators.  SNAKE couples state-space-adaptive screening to candidate--position records, reduced-window Schur reranking, beam exploration, recoverability-based generator ablation, and mixed fermion--boson validation.  To refer to the algorithm, we use the backronym SNAKE: Selection--Novel ADAPT Kost Evaluator, which emphasizes the novelty $\mathcal N$ and cost $K$ terms as well as the generator string's head growth and tail truncation as it moves through state space.  Phase I funnels the $O(\sum_{m\in\Gset}|\Pcal_m|)$ record surface through lower-confidence Fubini--Study trust-region gain divided by compiled and measurement cost $K(r;t)$.  Phase II adds tangent novelty and curvature so directions already expressible by the current local span are downweighted.  Phase III performs reduced-window Schur reranking after local refit relaxation, directly testing whether the candidate keeps value after nearby coordinates respond.  Optional batching and beam continuation handle compatible near-ties and early ordering ambiguity, while rollback-safe generator ablation removes verified stale generators.  Operationally, SNAKE targets the accuracy--resource surface directly: at fixed accuracy, it selects adaptive ansätze with fewer two-qubit gates, smaller depth, fewer effective parameters, or lower shot cost than fixed ansatz VQE and regular append-only ADAPT.

\section{SNAKE ADAPT ansatz-construction procedure}
\label{sec:static_method}

\begin{widetext}
\begin{center}
\includegraphics[width=\textwidth]{paper_I_static_adapt_root_images/static_adapt_selector_flow_attatched.png}
\captionof{figure}{SNAKE ADAPT ansatz-construction flow. Candidate--position records are generated from the operator pool, first filtered by first order energy gradient and hardware cost penalties; operator-commutation-class shortlists are further filtered by tangent span novelty and curvature; finally a Schur gain phase 3 scores the shortlist, and batch-compatible top scores are admitted where universe-branching fires if scores are degenerate. Cached history then probes for archaic-utility generators to remove prior to repeating the selection step.}
\label{fig:static_adapt_selector_flow}
\end{center}
\end{widetext}

\subsection{Three-phase generator-position ranking}

Contrary to traditional ADAPT, which appends a generator $A$ at the tail of the gate-sequence ansatz, SNAKE considers a live window of available insertion points within the existing gate sequence $U(\theta)$. The effect of inserting generator $A$ at position $p$ is equivalent to conjugation of $A$ by the subsequent ansatz gates, which is equal to the partial derivative with respect to the inserted parameter $\alpha$. That is, $[U(\theta)]_r^+|\psi \rangle=U(\theta)\oplus_p A_m |\psi \rangle= U_{>p}A_mU_{>p}^{\dagger}|\psi \rangle=\partial_\alpha |\psi_r(\theta \oplus \alpha)\rangle$. We denote these candidate--position record pairs as $r=(m,p)$, with admitted generator $m=A_m$. In the limit $p=|U(\theta)|+1$, we recover traditional append-only ADAPT.

Because subsequent candidate scoring mechanisms can be expensive, we introduce a tripartite scoring mechanism that incrementally reduces the candidate set while increasing the measurement burden. Briefly, Phase I scores the broad candidate--position surface with cheap uncertainty, hardware proxies, and a Fubini--Study trust region; algebraic shortlisting preserves support- and commutation-distinct survivors; Phase II reranks the retained records using inherited-window tangent novelty and curvature information; and Phase III applies the local reduced-window relaxation geometry used for final admission, batching, and beam continuation. Because shot cost increases with the expanded circuit, shortlisting pressure is iteration-dependent, and expensive phases may retire in the order Phase III then Phase II before termination by energy plateau or circuit-budget saturation.

\subsubsection{Phase I}
Phase 1 scoring takes the form of
\begin{equation}
S_1(r;t)=\frac{\Delta E_{1}(r)}{1+K(r;t)}.
\label{eq:phase1_score}
\end{equation}
where $K(r;t)$ denotes the hardware cost of candidate $r$, and $\Delta E_{1}(r)$ represents a shot-noise uncertainty-aware energy gradient within a Fubini--Study metric trust region \cite{Nocedal2006TrustRegion}. Explicitly,
\begin{equation}
\begin{split}
\Delta E_{1}(r)={}&\max_{|\alpha|\le \frac{\rho}{\sqrt{F(r)}}} \left[g(r)|\alpha|-\frac{\lambda_FF_{\rm raw}(r)\alpha^2}{2}\right].
\end{split}
\label{eq:phase1_gain}
\end{equation}
Where the energy gradient estimate,
\begin{equation}
g(r)=\left|2\Re\langle\psi|H|U_{>p}A_mU_{>p}^{\dagger}\ket{\psi}\rangle\right|-z_{\alpha}\sigma_r \geq0,
\label{eq:lcb_gradient}
\end{equation}
is computed within a confidence interval of standard error $\sigma_r$, and $z_{\alpha}$ is the one-sided normal quantile fixing the lower-confidence level \cite{Ikhtiarudin2025ShotEfficient}. The regularized raw tangent scale $F(r)=\langle\bar\tau_r,\bar\tau_r\rangle+\eps_F$ defines a single-coordinate trust-region proxy. Here $\bar \tau = Q_\psi\tau= Q_\psi\partial_\theta |\psi\rangle=(I-|\psi\rangle \langle \psi |) \partial_\theta |\psi \rangle$ is the tangent direction induced from $r$'s parameter $\theta$, projected onto the state's horizontal subspace. The Fubini--Study metric regularizes generator comparison by state displacement, whereas traditional ADAPT scores by unit parameter displacement and can bias local parameter speed rather than physical state-space motion \cite{ProvostVallee1980,Stokes2020QNG}.

The hardware-cost penalty $K(r;t)$, results from proxy rules and inexpensive pool metadata. It sums estimated two-qubit count, layout-sensitive two-qubit depth, new parameter cost, measurement-group or shot cost, and remaining-budget pressure; Here each cost-component is centered and scaled over the live record family before weighting such that candidate cost competition is closed. Measurement-grouping work supports treating Pauli burden as part of the adaptive algorithmic cost, rather than as post hoc accounting \cite{Verteletskyi2020CliqueCover,Yen2020CompatibleMeasurements}.

Consider the compiled generator, $A_m=\sum_{\ell}a_{m\ell}P_{m\ell}$ where $P$ is a tensor-product of Pauli-matricies tensor-plus identitity, $\sigma \in \{I,X,Y,Z\}$. Define the Pauli-support, $S:=\operatorname{supp}(P)=\{q:P_q\ne I\}$. Then, for a Pauli exponential compiled by a CNOT ladder, the total 2q gate count is given by,
\begin{equation}
\widehat C_{\twq}(r)=\sum_{\ell}2\,[|S_{m\ell}|-1]_+,
\label{eq:pauli_2count_proxy}
\end{equation}
whereas the layout-aware two-qubit depth is given by,
\begin{equation}
\widehat D_{\twq}(r)=\sum_{\ell}2\,\operatorname{span}_{\mathcal L}
\bigl(S\bigr),
\label{eq:pauli_2depth_proxy}
\end{equation}
where $\operatorname{span}_{\mathcal L}(S)$ is the number of edges in the shortest connected subgraph of the qubit-hardware QPU graph, $\mathcal L = (V,E)$ containing $S$, and is zero for $|S|\le1$. The parameter proxy is the precomputed number of new free amplitudes for a given generator. For measurement burden, the gradient observable $O_r=i[H,\widetilde A_r]$ is partitioned into qubit-wise commuting groups $O_{r,b}=\sum_{\ell\in B_{r,b}}c_{r,b\ell}P_{r,b\ell}$. If $\mathcal B_r^{\rm new}$ denotes groups not already covered by measurement reuse, the shot proxy scales as,
\begin{equation}
\widehat C_{\rm shot}(r)\propto
\frac{1}{\sigma_\star^2}
\left(
\sum_{b\in\mathcal B_r^{\rm new}}
\sqrt{\sum_{\ell\in B_{r,b}}|c_{r,b\ell}|^2}
\right)^2,
\label{eq:shot_cost_proxy}
\end{equation}
with $\sigma_\star$ the target standard error for the candidate-gradient estimate \cite{Ikhtiarudin2025ShotEfficient}. Final reported resources are always obtained from the frozen compiler and measurement allocator; these proxies only rank live candidates. Candidates of a sufficiently high score are then shortlisted by algebraic type before Phase II.

\subsubsection{Algebraic shortlisting}
After Phase I assigns each candidate--position record a trust-region and cost score, SNAKE preserves algebraically distinct survivors before the more expensive tangent-geometry evaluation.  For a candidate record $r=(m,p)$, let $\mathcal C_1(r;t)$ denote the generators in its local candidate context, including local pool neighbors and inherited-window generators.  For each $A_n\in\mathcal C_1(r;t)$, the selector checks support overlap and commutation,
\begin{equation}
\operatorname{supp}(A_m)\cap\operatorname{supp}(A_n),
\qquad
[A_m,A_n].
\label{eq:algebraic_shortlist_checks}
\end{equation}
For Pauli-expanded generators, let $N_{\rm ac}(m,n)$ be the number of sites in the shared support where $P_q^{(m)}$ and $P_q^{(n)}$ anticommute.  The commutation test is
\begin{equation}
[A_m,A_n]=0
\quad\Longleftrightarrow\quad
N_{\rm ac}(m,n)\equiv 0 \pmod 2.
\label{eq:pauli_commutation_parity}
\end{equation}
The resulting four-way shortlist preserves support-overlapping commuting records, support-overlapping noncommuting records, support-disjoint commuting records, and mixed or unresolved records.  Within each type, records remain ordered by the Phase-I trust-region and cost score, so this step preserves algebraically distinct candidates before tangent geometry is measured.

\subsubsection{Phase II}
Typical ADAPT's ignorance of ansatz-space expressibility can admit operators that induce state shifts already expressible by simple coordinate refit. As such, introducing a reward for geometrically novel state-space directions improves convergence to lower-energy adaptive ansätze. We denote the $[0,1]$-ranged geometric novelty reward as $\mathcal N(r;t)$.

The phase II score then becomes,
\begin{equation}
S_2(r;t)=\frac{\Delta E_{2}(r)\,\mathcal N_2(r;t)}{1+K_2(r;t)}.
\label{eq:phase2_score}
\end{equation}
where
\begin{equation*}
\mathcal N_2(r;t)
=
\operatorname{clip}_{[0,1]}
\left(
1-
\frac{
q_r^{\top}(Q_r+\lambda_{\rm nov}I)^{-1}q_r
}{
F_{\rm raw}(r)
}
\right),
\end{equation*}
with
\begin{equation*}
(Q_r)_{ij}=\Re\langle\tau_i,\tau_j\rangle,
\qquad
(q_r)_i=\Re\langle\tau_i,\tau_r\rangle,
\qquad
i,j\in W_r.
\end{equation*}
Here $W_r$ is an inherited overlap window, $Q_r$ is its Fubini--Study Gram block, and $q_r$ is the candidate-to-window tangent covariance vector.  The factor $\mathcal N_2(r;t)$ is therefore a regularized residual-span score: it is small when $\tau_r$ is already expressible by the inherited tangent family and large when the candidate contributes a new state-space direction.

We then use second order energy estimation
\begin{equation*}
\Delta E_{2}(r)
=\max_{|\alpha|\le \frac{ \rho}{\sqrt{F(r)}}}
\Bigl[g(r)|\alpha| -\frac{h_r\alpha^2}{2}\Bigr],
\end{equation*}
where the non-negative $h_r=2\Re\!\left( ( \partial_{\alpha}\langle \psi|)H|\partial_{\alpha}\psi\rangle+\langle\psi|H|\partial^2_{\alpha}\psi\rangle\right)$, and $g(r)$ is the cashed gradient information from phase I. Here the second derivative of the current state is found from $\partial^2_{\alpha}|\psi \rangle=-(U_{>p}A_mU_{>p}^{\dagger})^2|\psi\rangle$.

\subsubsection{Phase III}
Phase 3 performs reduced-window reranking.  This is the optional higher-cost Schur extension of the core Phase-I/II selector: if it is disabled, SNAKE reduces to geometry- and cost-screened ADAPT with ordinary local refit. The candidate is inserted locally, a reduced refit window $W_r$ is opened, and the local quadratic model
\begin{equation}
\begin{split}
E(\alpha,\delta\theta_{W_r};r)\approx{}&E_0+g_r\alpha+\frac12h_r\alpha^2 \\
&+\alpha b_r^{\top}\delta\theta_{W_r} +\frac12\delta\theta_{W_r}^{\top}H_{W_r}\delta\theta_{W_r}
\end{split}
\label{eq:local_quadratic}
\end{equation}
is used, where we eliminate window variables $\delta \theta_W$ to produce a Schur-complement 1-D model based on candidate parameter $\alpha$.
The treatment is as follows:
Let $b_r$ and $H_{W_r}$ be the mixed vector and window Hessian block respectively of the local model, defined by $b_j:=
\left.
\frac{\partial^2 E(\alpha,\delta\theta_W)}
{\partial\alpha\,\partial(\delta\theta_j)}
\right|_{\alpha=0,\delta\theta_W=0}$, and $(H_W)_{jk}:=
\left.
\frac{\partial^2 E(\alpha,\delta\theta_W)}
{\partial(\delta\theta_j)\,\partial(\delta\theta_k)}
\right|_{\alpha=0,\delta\theta_W=0},$ for $j,k\in W_r.$
For the expected energy benefit after window refit, SNAKE predicts the refit response after generator insertion. The required quantities are the Schur-complemented relaxed curvature of candidate $r$, $(Q_r)_{jk}=\Re\langle\bar\tau_j,\bar\tau_k\rangle,$ the window-tangent Gram matrix $(q_r)_j=\Re\langle\bar\tau_j,\bar\tau_r\rangle$, and the noise-regularized window Hessian $\mathbf R_r=\frac12(H_{W_r}+H_{W_r}^{\top})+\lambda_H I$, symmetrized and shifted to ensure positive definiteness.

Thus the regularized Schur complement relaxed curvature becomes, $\widetilde h_r=h_r-b_r^{\top}\mathbf R_r^{-1}b_r,$ the Fubini-Study norm of the candidate after subtracting state-space displacement due to window relaxation becomes $F_{\rm red}=F(r)-2q_r^{\top}\mathbf R_r^{-1}b_r+b_r^{\top}\mathbf R_r^{-1}Q_r\mathbf R_r^{-1}b_r,$ the relaxed overlap after refit becomes $q_r^{\rm red}=q_r-Q_r\mathbf R_r^{-1}b_r$, which motivates the geometrical novelty factor post-window refit:
\begin{equation}
\mathcal N_3(r;t)=1-\frac{(q_r^{\rm red})^{\top}(Q_r+\eps_{\rm nov}I)^{-1}q_r^{\rm red}}
{F_{\rm red}}
\qquad \in [0,1] .
\end{equation}
Then,
\begin{equation}
\Delta E_{\TR}(r)=
\max_{|\alpha|\le \frac{\rho}{\sqrt{F_{\rm red}(r)}}}
\left[g(r)|\alpha|-\frac{\widetilde h_r \alpha^2}{2}\right].
\label{eq:phase3_gain}
\end{equation}
And consequently the authoritative single-record score is
\begin{equation}
\begin{split}
S_3(r;t)=\frac{\Delta E_{3}(r)\,\mathcal N_3(r;t)}{K_3(r;t)}
\end{split}
\label{eq:phase3_score}
\end{equation}

\subsubsection{Batching and Beaming}
The phase III results produces a ranked scoring set, $\mathcal R_3=\{(r,p) \mid r \in \mathcal A, p\in W, S_3(r,p) \geq \tau\}$. Batching is an optional admission accelerator layered on the same candidate--position score. We employ a TETRIS-ADAPT-VQE-inspired batching criterion; TETRIS-ADAPT-VQE appends batched operators under mutual-commutation, disjoint-qubit action, and gradient-threshold conditions \cite{Anastasiou2024}. We use these conditions, together with joint insertability, support overlap, layout scheduling, and measurement-reuse checks, as a shortlisting mechanism before a more shot-intensive geometric distinguishability test for reduced tangent independence/additivity.

Consider further degenerately scored batched-admission candidates. Beam continuation is likewise optional: disabling it leaves a single greedy branch closer to ordinary ADAPT, while enabling it preserves several high-value early operator-order hypotheses until the local evidence separates them. A beam treatment splits into $k$ parallel ADAPT quantum circuits, each run for $i$ iterations until a winning circuit is discernible, after which losing branches are terminated. SNAKE employs this same treatment with the three-phase selection mechanism and a specific iteration-index dependency described in the next section.
% \placeholder{Citation needed for beam/branch continuation; identify a direct prior-work anchor if one is retained.}
% \placeholder{Add candidate-record deduplication by ansatz identity and rounded amplitudes if retained.}

\subsection{Iteration-Dependent Gate-Selection \& Other Features}
Although the tripartite scoring sequence evaluation occurs at a fixed iteration indicie, the weightings of its features, such as novelty reward, energy drop, and cost have monotonic time-dependency. Their weighted behavior is determined by a useful-runway forecast estimator which determines how many profitable insertions remain under the live depth and shot budgets, with further possible bounds being the signal to noise-ratio and energy-plateau behavior.

Specifically, novelty and cost are respectively rewarded and penalized more early, whereas energy drop is rewarded more later. Furthermore, insertion position candidates, refit windows, and short-listing sizes are down-sized with iterations in order to combat the increasing shot cost from an expanded circuit, and to counter-balance a mechanism which increases the shot-count of candidates in late iteration-horizon scoring such that the typically decreasing signal-to-noise ratio behavior is nullified. Such shot-count is further combatted by phase retirement -- namely, the scoring pipeline retires expensive phases in the order Phase 3 then Phase 2, with eventual algorithm termination according to either circuit-budget saturation, or vanishing projected gradients/plateaus of energy drops satisfying a patience interval.

Increased shortlist pressure indirectly affects batching, whereas beam continuation is biased early because circuit splitting is cheaper and early operator structure has higher leverage. These controls are not separate algorithms; they are limiting-case switches inside the same SNAKE selection rule.

Although the emphasis on geometrical utility increases the probability of late-stage generator usefulness, certain generators can become structurally redundant after subsequent admissions and refits. Formally, coordinate $j$ is redundant at tolerance $\epsilon_{\rm pr}$ if deleting its generator admits a constrained refit whose optimized energy regression is at most $\epsilon_{\rm pr}$; the generator-ablation procedure below approximates this recoverability condition before committing an ablation. As such, a tracking mechanism of an operator's proxied utility is employed, and we bias candidates that, after satisfying an admission-recency cool-down, exhibit historic large-amplitude initialization followed by collapse to a stably small amplitude due to window refit. A local quadratic Taylor model using cached information then approximates post-ablation refit-window compensation through the Schur complement.

See appendix~\ref{app:static_controller_details} for explicit telemetry, SNR-floor, shortlist, and remaining useful runway formulas.


\subsection{Encoding and layout surface}
\label{sec:encoding_surface}

SNAKE is defined over mapped Pauli generators, but the resource consequences depend on the mapping that produced them. We therefore treat the encoding as part of the benchmark surface rather than as a hidden preprocessing detail. Fermion mappings include Jordan--Wigner, parity, and Bravyi--Kitaev; boson encodings include binary, Gray-code, unary, and variational basis-state registers; and register layouts include spin-blocked fermions, site-interleaved fermions, and site-local fermion--boson blocks. Specific phonon-boson qubit encodings have been studied in digital and hybrid quantum-simulation settings \cite{Macridin2018,Sawaya2020,Li2023,Denner2023}, and these encodings are made explicit in Appendix~\ref{app:encodings}. The selector is not allowed to change the physics Hamiltonian across this surface, but it may observe the compiled support, two-qubit-depth proxy, and measurement grouping induced by each encoding. The canonical formulas are collected in Appendix~\ref{app:encodings}, while Appendix~\ref{app:evaluation_details} records the reduced encoding panel used to separate physics accuracy from mapping-dependent resource variance.


\section{Results}
\label{sec:results}
\label{sec:benchmarks}

The results are organized around accuracy--resource Pareto comparisons. We compare final ground-state energy accuracy, state fidelity, compiled two-qubit count, compiled depth, and shot-cost proxy against fixed ansatz VQE, append-only ADAPT, Qubit-ADAPT-VQE, TETRIS-ADAPT-VQE, and Geo-ADAPT-VQE. We then isolate the contribution of the three-phase SNAKE selector, hardware-cost terms, tangent novelty, local-window refit, batching, beam continuation, and generator ablation through paired ablations. The Hubbard--Holstein rows provide, to our knowledge, the first ADAPT-VQE validation on mixed fermion--boson Hamiltonians rather than purely fermionic molecular Hamiltonians. Each reported cell is tied to the corresponding benchmark observable rather than inferred from a different metric.

\subsection{Ground-state results}

\structlabel{Ground-state results}Table~\ref{tab:static_claims} gives the ground-state comparison by Hamiltonian class averages. The same-cutoff columns test whether the adaptive ansatz can reproduce the ED target in the working Hilbert space \cite{Weisse2008}; the $n_{\rm ph}=4$ columns test the residual cutoff error for phonon-bearing systems after embedding the lower-cutoff variational state. The fixed VQE ansatz is class matched: UCCSD for fermions; displacement, squeezing/quadrature, hopping, and number/Kerr blocks for bosons; and lifted UCCSD plus Lang--Firsov, displacement/quadrature, and electron--phonon dressing blocks for mixed systems. Append-only ADAPT uses Hamiltonian Pauli-term pools; Qubit-ADAPT-VQE, TETRIS-ADAPT-VQE, and Geo-ADAPT-VQE use pairwise qubit-excitation pools with their respective selection rules; SNAKE uses the corresponding class-matched geometry--cost pool. Resource columns report final compiled two-qubit count, two-qubit depth when available, total compiled depth, and total shot-cost proxy. The fermion--boson aggregate uses Hubbard--Holstein and spin--boson benchmark rows; harmonic--Kerr stress tests are excluded from that aggregate because they do not meet the convergence gate.

\begin{widetext}
\inlinetablecaption{tab:static_claims}{Ground-state Pareto benchmark by Hamiltonian class and method.}
\begin{center}
\scriptsize
\begin{ruledtabular}
\begin{tabular*}{\textwidth}{@{\extracolsep{\fill}}p{0.105\textwidth}p{0.195\textwidth}cccccccc@{}}
Class & Method & $|\Delta E|_{\rm same}$ & $|\Delta E|_{4}$ & $1-F_{\rm same}$ & $1-F_{4}$ & $N_{\twq}$ & $D_{\twq}$ & $D_{\rm circ}$ & $S_{\rm shot}$ \\
\colrule
fermionic & fixed ansatz VQE & \tentative{$0.3425$} & -- & -- & -- & \tentative{$8.2$} & -- & \tentative{$36$} & \tentative{$5.34\!\times\!10^{7}$} \\
fermionic & append-only ADAPT & \tentative{$3.17$} & -- & -- & -- & \tentative{$0$} & -- & \tentative{$5$} & \tentative{$1.33\!\times\!10^{5}$} \\
fermionic & Qubit-ADAPT-VQE & \tentative{$0.2234$} & -- & -- & -- & \tentative{$6$} & -- & \tentative{$9$} & \tentative{$1.11\!\times\!10^{6}$} \\
fermionic & TETRIS-ADAPT-VQE & \tentative{$0.2322$} & -- & -- & -- & \tentative{$6$} & -- & \tentative{$9$} & \tentative{$1.30\!\times\!10^{6}$} \\
fermionic & Geo-ADAPT-VQE & \tentative{$0.2234$} & -- & -- & -- & \tentative{$6$} & -- & \tentative{$9$} & \tentative{$1.24\!\times\!10^{6}$} \\
fermionic & SNAKE & \tentative{$5.74\!\times\!10^{-8}$} & -- & \placeholder{...} & -- & \tentative{$51.1$} & \placeholder{...} & \tentative{$199.4$} & \tentative{$875.9$} \\
\colrule
bosonic & fixed ansatz VQE & \tentative{$0.0803$} & -- & -- & -- & \tentative{$8$} & -- & \tentative{$36$} & \tentative{$3.97\!\times\!10^{7}$} \\
bosonic & append-only ADAPT & \tentative{$0.9501$} & -- & -- & -- & \tentative{$0$} & -- & \tentative{$2.5$} & \tentative{$1.23\!\times\!10^{6}$} \\
bosonic & Qubit-ADAPT-VQE & \tentative{$0.5597$} & -- & -- & -- & \tentative{$2$} & -- & \tentative{$3$} & \tentative{$6.72\!\times\!10^{5}$} \\
bosonic & TETRIS-ADAPT-VQE & \tentative{$0.5597$} & -- & -- & -- & \tentative{$2$} & -- & \tentative{$3$} & \tentative{$6.72\!\times\!10^{5}$} \\
bosonic & Geo-ADAPT-VQE & \tentative{$0.5597$} & -- & -- & -- & \tentative{$2$} & -- & \tentative{$3$} & \tentative{$9.78\!\times\!10^{5}$} \\
bosonic & SNAKE & \tentative{$1.85\!\times\!10^{-7}$} & \placeholder{...} & \placeholder{...} & \placeholder{...} & \tentative{$106.5$} & \placeholder{...} & \tentative{$387.5$} & \placeholder{...} \\
\colrule
fermion--boson & fixed ansatz VQE & \tentative{$0.2345$} & -- & -- & -- & \tentative{$11.6$} & -- & \tentative{$18.6$} & \tentative{$3.29\!\times\!10^{7}$} \\
fermion--boson & append-only ADAPT & \tentative{$2.46$} & -- & -- & -- & \tentative{$5.5$} & -- & \tentative{$38.5$} & \tentative{$6.34\!\times\!10^{5}$} \\
fermion--boson & Qubit-ADAPT-VQE & \tentative{$0.2147$} & -- & -- & -- & \tentative{$4$} & -- & \tentative{$6$} & \tentative{$4.73\!\times\!10^{6}$} \\
fermion--boson & TETRIS-ADAPT-VQE & \tentative{$0.3310$} & -- & -- & -- & \tentative{$4$} & -- & \tentative{$6$} & \tentative{$2.50\!\times\!10^{6}$} \\
fermion--boson & Geo-ADAPT-VQE & \tentative{$0.2147$} & -- & -- & -- & \tentative{$4$} & -- & \tentative{$6$} & \tentative{$6.10\!\times\!10^{6}$} \\
fermion--boson & SNAKE & \tentative{$4.89\!\times\!10^{-5}$} & \placeholder{...} & -- & \placeholder{...} & \tentative{$69$} & \placeholder{...} & \tentative{$163$} & -- \\
\colrule
all averaged & fixed ansatz VQE & \tentative{$0.2668$} & -- & -- & -- & \tentative{$9.0$} & -- & \tentative{$31.9$} & \tentative{$4.59\!\times\!10^{7}$} \\
all averaged & append-only ADAPT & \tentative{$2.58$} & -- & -- & -- & \tentative{$1.1$} & -- & \tentative{$11.2$} & \tentative{$4.54\!\times\!10^{5}$} \\
all averaged & Qubit-ADAPT-VQE & \tentative{$0.2889$} & -- & -- & -- & \tentative{$4.8$} & -- & \tentative{$7.2$} & \tentative{$1.75\!\times\!10^{6}$} \\
all averaged & TETRIS-ADAPT-VQE & \tentative{$0.3175$} & -- & -- & -- & \tentative{$4.8$} & -- & \tentative{$7.2$} & \tentative{$1.42\!\times\!10^{6}$} \\
all averaged & Geo-ADAPT-VQE & \tentative{$0.2889$} & -- & -- & -- & \tentative{$4.8$} & -- & \tentative{$7.2$} & \tentative{$2.16\!\times\!10^{6}$} \\
all averaged & SNAKE & \tentative{$7.01\!\times\!10^{-5}$} & \placeholder{...} & \tentative{$2.98\!\times\!10^{-5}$} & \placeholder{...} & \tentative{$78$} & \tentative{$243$} & \placeholder{...} & \tentative{$4.0\!\times\!10^{3}$} \\
\end{tabular*}
\end{ruledtabular}
\end{center}
\end{widetext}

Across the fermion--boson block, the lowest-error non-SNAKE competitors are Qubit-ADAPT-VQE and Geo-ADAPT-VQE, both at \tentative{$0.2147$} mean energy error.  SNAKE reduces this error to \tentative{$4.89\!\times\!10^{-5}$}, a \tentative{$99.98\%$} reduction.  The corresponding compiled-resource columns show an accuracy--resource tradeoff rather than uniform resource dominance: SNAKE uses \tentative{$69$} compiled two-qubit gates and \tentative{$163$} compiled depth, compared with \tentative{$4$} and \tentative{$6$} for those two low-error qubit-excitation baselines.

Table~\ref{tab:ablation_matrix} isolates SNAKE mechanism efficacy by disabled-minus-full comparison. Hamiltonian, seed, cutoff, encoding, backend, and optimizer settings are fixed. Positive $\delta|\Delta E|$ means worse final energy accuracy; positive resource deltas mean heavier compilation or overhead cost. The shot column is the live selector proxy: nominal selector-measurement work spent by shortlist, rerank, beam, batch, generator-ablation, and rollback decisions. It excludes final Hamiltonian readout shots and calibrated hardware-shot accounting. Daggered rows are targeted stress ablations.

\begin{widetext}
\inlinetablecaption{tab:ablation_matrix}{SNAKE ablations. Entries are mean paired changes from ablation runs. $\delta|\Delta E|$ measures the change in energy-error from introducing the ablation, and $N_{\twq}$ is compiled two-qubit gate count, $D_{\twq}$ is compiled two-qubit depth, $D_{\rm circ}$ is compiled circuit depth, and $S_{\rm ctrl}^{\rm proxy}$ is SNAKE nominal selector-shot work proxy. Positive results indicate worse accuracy or added overhead.}
\begin{center}
\begin{ruledtabular}
\begin{tabular}{p{0.33\textwidth}ccccc}
Variant & mean $\delta|\Delta E|$& mean $\Delta N_{\twq}$ & mean $\Delta D_{\twq}$ & mean $\Delta D_{\rm circ}$ & mean $\Delta S_{\rm ctrl}^{\rm proxy}$ \\
\colrule
no hardware cost $K$ & $-4.92\!\times\!10^{-8}$ & $+30.3$ & \placeholder{...} & $+122.3$ & $0$ \\
no tangent novelty factors $(\mathcal N_2,\mathcal N_3)$ & $+2.64\!\times\!10^{-4}$ & $+21.3$ & \placeholder{...} & $+97.7$ & $0$ \\
no Phase-2 tangent novelty $\mathcal N_2$ & $-4.92\!\times\!10^{-8}$ & $+30.3$ & \placeholder{...} & $+122.3$ & $0$ \\
no Phase-3 reduced-window geometry factor $\mathcal N_3/\Delta^{\rm red}$ & $+2.64\!\times\!10^{-4}$ & $-9.3$ & \placeholder{...} & $-26.0$ & $0$ \\
Phase-1+2 only; no Phase 3 & $+9.43\!\times\!10^{-6}$ & $-27.7$ & \placeholder{...} & $-117.0$ & $0$ \\
Phase-1 only; no Phase 2/3 & $+9.61\!\times\!10^{-5}$ & $-52.7$ & \placeholder{...} & $-139.3$ & $0$ \\
append-only refit; no active local window& $+6.33\!\times\!10^{-2}$ & $-37.0$ & \placeholder{...} & $-142.0$ & $0$ \\
no generator ablation& $+2.39\!\times\!10^{-5}$ & $+60.7$ & \placeholder{...} & $+220.3$ & $+770.0$ \\
no batching & $+1.65\!\times\!10^{-4}$ & $+2.0$ & \placeholder{...} & $+30.0$ & $0$ \\
no beam continuation & $+9.43\!\times\!10^{-5}$ & $-65.7$ & \placeholder{...} & $-216.0$ & $0$ \\
regular append-only ADAPT & $+2.35\!\times\!10^{-2}$ & $-48.3$ & \placeholder{...} & $-148.7$ & $0$ \\
\end{tabular}
\end{ruledtabular}
\end{center}
\end{widetext}

The ablation rows separate three effects. Cost and novelty terms mostly change which operator sequence is purchased for a fixed selector schedule. Local-window refit is accuracy-positive in the stressed Hubbard--Holstein and extended-Hubbard points, while generator ablation trades an explicit recoverability test for lower stale-depth burden. Batching, beam continuation, Schur reranking, and generator ablation are optional extensions of the core selector: the corresponding ablation rows show limiting cases closer to one-record or append-only ADAPT procedures, not unrelated algorithms. The phase-retirement controls show the hierarchy: retiring Phase 3, and then Phase 2, reduces immediate cost but also removes the geometry tests that protect late-stage admissions.


\section{Discussion and conclusion}
\label{sec:discussion}
\label{sec:conclusion}

\structlabel{Mechanistic interpretation}
SNAKE's mechanism is selective acquisition of useful tangent directions under finite circuit and measurement budgets. Candidate--position records expose insertion order, Fubini--Study scaling compares state displacement rather than raw parameter displacement, and Schur-complement reranking asks whether a candidate remains useful after nearby coordinates refit. Cost penalties, batching, beam branching, and generator ablation then keep this geometry from becoming a depth-only construction.

\structlabel{Scope and limitations}
The study is a compiled, finite-shot-accounted small-system benchmark with ED used as the reference. Its scope is ADAPT ansatz construction: class-wise Hamiltonian sweeps, fidelity extraction, high-cutoff phonon stress, calibrated noise, mitigation, and hardware validation enter as resource and validation axes for the same SNAKE acquisition rule. The encoding and ordering panel remains part of the method evaluation because SNAKE's hardware-cost terms make encoding choice part of the resource claim.

\structlabel{Summary}
We presented SNAKE, a geometry- and cost-aware ADAPT procedure for mixed fermion--boson Hamiltonians. SNAKE ranks insertion candidates by trust-region gain, tangent novelty, reduced-window Schur geometry, and compiled/measurement cost, while using batching, beam continuation, and rollback-safe generator ablation to maintain a compact adaptive ansatz. The mixed-system rows quantify the energy and compiled-cost comparison, and the class-stratified table organizes the corresponding fidelity, cutoff-stress, two-qubit-depth, shot-cost, and all-class axes under the same Pareto protocol.


\clearpage
\appendix


\section{Tentative table provenance ledger}
\label{app:data_ledger}

\tentative{Draft-only provenance appendix for Table~\ref{tab:static_claims} and Table~\ref{tab:ablation_matrix}. This section is for repo-agent maintenance and must be removed, hidden, or converted into archival support material before journal submission.}

\subsection{Table I source map}
\label{app:table_i_source_map}

Table~\ref{tab:static_claims} combines recovered non-SNAKE competitor rows with SNAKE rows. The non-SNAKE rows come from the recovered CHTC table suite; the SNAKE rows are separate SNAKE-support artifacts and should not be overwritten by the competitor aggregator.

\begin{itemize}
\item \textbf{Recovered non-SNAKE competitor suite.} Records file: \path{chtc/phase3_optuna/input/generic_static_table_records_recovered_no_h2o_hea.tsv}. Raw output root: \path{raw_outputs/chtc_phase3_optuna/generic_static_table/}. Validation summary: \path{raw_outputs/generic_static_table_full_output_check_recovered_no_h2o_hea.json}. Current validation status: expected $61$, benchmarked $61$, missing $0$, unusable $0$, quality-nonpassing $0$.
\item \textbf{Recovered-row aggregator.} Regenerate the non-SNAKE rows with \path{pipelines/exact_bench/summarize_table_i_static_results.py}. The current generated outputs are \path{raw_outputs/table_i_static_paper_recovered_no_h2o_hea/table_i_static_results_summary.json}, \path{raw_outputs/table_i_static_paper_recovered_no_h2o_hea/table_i_static_results_rows.csv}, and \path{raw_outputs/table_i_static_paper_recovered_no_h2o_hea/table_i_static_claim_rows.tex}.
\item \textbf{Paper-label mapping for recovered rows.} \path{static_hea_qiskit_vqe} maps to fixed ansatz VQE; \path{static_qiskit_adapt_vqe} maps to append-only ADAPT; \path{static_qubit_qeb_adapt_vqe} maps to Qubit-ADAPT-VQE; \path{static_tetris_qubit_adapt_vqe} maps to TETRIS-ADAPT-VQE; \path{static_geo_qubit_adapt_vqe} maps to Geo-ADAPT-VQE. The historical source key containing \path{qeb} is an internal implementation key only; the manuscript label is Qubit-ADAPT-VQE.
\item \textbf{Fermionic SNAKE row.} Remote CHTC root: \path{/home/jsstrobel/Holstein_phase3_optuna_chtc}. The current collective-span record is \path{tripartite_spsa_A_current_collective_fermionic_smallrobust_target5e5_v1}, CHTC job $6305147.0$, with trace targets \path{raw_outputs/tripartite_spsa_A_current_collective_fermionic_smallrobust_target5e5_v1/summary.json}, \path{raw_outputs/tripartite_spsa_A_current_collective_fermionic_smallrobust_target5e5_v1/command.sh}, and \path{logs/phase3-tripartite-spsa-ab.6305147.0.*}. The historical comparison record is \path{tripartite_spsa_B_pre_phase2_novelty_legacy_fermionic_smallrobust_target5e5_v1}, CHTC job $6305147.3$.
\item \textbf{Bosonic and fermion--boson SNAKE rows.} The L2/$n_{\rm ph}=1$ reset manifest is \path{chtc/phase3_optuna/input/global_tripartite_spsa_ab_l2nph1_reset_manifest.json}, generated by \path{chtc/phase3_optuna/generate_tripartite_spsa_ab_l2nph1_reset_records.py} from \path{chtc/phase3_optuna/input/global_tripartite_spsa_ab_records.tsv}. Current reset output roots include \path{raw_outputs/reset6305147_spsa_A_bosonic_bose_hubbard_L2_nph1_smallrobust_target5e5_v1} and \path{raw_outputs/reset6305147_spsa_A_mixed_hh_L2_nph1_smallrobust_target5e5_v1}. If a future Table~\ref{tab:static_claims} fermion--boson aggregate uses both Hubbard--Holstein and spin--boson SNAKE rows, add the spin--boson SNAKE artifact path here before promoting the aggregate.
\end{itemize}

\subsection{Table II source map}
\label{app:table_ii_source_map}

Table~\ref{tab:ablation_matrix} is currently a local-pilot paired-delta table, not the final canonical ablation sweep. A repo agent updating this table should replace the whole row block only after locating or regenerating normalized machine-readable full-versus-ablated artifacts.

\begin{itemize}
\item \textbf{Current machine-readable table aggregate.} The values currently carried into Table~\ref{tab:ablation_matrix} trace to \path{prompt-exports/phase3_static_adapt_ablation_matrix_native_sctrl_v2.json}. That file aggregates \path{prompt-exports/phase3_static_adapt_ablation_matrix_v1.json}, \path{tmp/static_phase_route_ablation_v1/phase_route_ablation_results.json}, \path{tmp/static_prune_pressure_sweep_v1/prune_pressure_results.json}, and \path{tmp/static_window_refit_sweep_v1/window_refit_results.json}.
\item \textbf{Human-readable local-pilot notes.} The local-pilot support notes are \path{prompt-exports/phase3_static_adapt_ablation_matrix_v1.md}, \path{prompt-exports/2026-04-30-static-phase-route-ablation-v1.md}, \path{prompt-exports/2026-04-30-static-prune-pressure-sweep-v1.md}, and \path{prompt-exports/2026-04-30-static-window-refit-sweep-v1.md}.
\item \textbf{Sign convention.} Table~\ref{tab:ablation_matrix} uses disabled-minus-full deltas: positive energy means worse accuracy, and positive resource entries mean higher cost. The generator-ablation row is sign-flipped relative to the source row labeled \path{static pruning recoverability ladder vs no pruning} in \path{prompt-exports/phase3_static_adapt_ablation_matrix_native_sctrl_v2.json}; the manuscript row reports \emph{no generator ablation} as the disabled-feature case.
\item \textbf{Replacement contract.} Before replacing Table~\ref{tab:ablation_matrix}, follow \path{prompt-exports/2026-04-30-static-ablation-reproduction-agent-spec.md}: freeze one canonical policy, run full SNAKE and matched ablations at identical Hamiltonian, seed, cutoff, encoding, backend, and resource settings, emit paired-delta JSON and aggregate JSON, and keep selector-work proxy fields separate from calibrated physical shot counts.
\item \textbf{Related batch report.} The current high-level SNAKE batch report is \path{output/pdf/phase3_canonical_three_lane_spsa_vs_powell_20260508.pdf}, generated by \path{tmp/pdfs/canonical_batches/make_canonical_report.py}. Treat that PDF as context and status evidence, not as the primary source for replacing Table~\ref{tab:ablation_matrix} if JSON/CSV artifacts are available.
\end{itemize}

\section{Benchmark Hamiltonian families}
\label{app:hamiltonian_families}

The benchmark registry contains spinful Hubbard variants, Hubbard--Holstein instances, spinless $t$--$V$ chains, boson-only chains, spin-boson/generalized-Rabi controls, and restricted closed-shell molecular tests. Onsite coefficients below absorb the sign convention used by a particular benchmark specification; the operator content is the invariant part of the family definition.

The spinful lattice base is
\begin{align}
H_{\rm Hub}[\epsilon]={}&
-t\sum_{\langle i,j\rangle,\sigma}
\left(c_{i\sigma}^{\dagger}c_{j\sigma}
+c_{j\sigma}^{\dagger}c_{i\sigma}\right)
\nonumber\\
&
+U\sum_i n_{i\uparrow}n_{i\downarrow}
+\sum_{i,\sigma}\epsilon_i n_{i\sigma}.
\label{eq:app_hubbard}
\end{align}
The registered ionic, extended, and next-nearest-neighbor variants are
\begin{align}
H_{\rm ion}&=H_{\rm Hub}[\epsilon_i=(-1)^i\Delta],
\nonumber\\
H_{\rm ext}&=H_{\rm Hub}[\epsilon]+V\sum_{\langle i,j\rangle} n_i n_j,
\nonumber\\
H_{t,t',U}&=H_{\rm Hub}[\epsilon]
\nonumber\\
&\quad
-t'\sum_{\langle\!\langle i,j\rangle\!\rangle,\sigma}
\left(c_{i\sigma}^{\dagger}c_{j\sigma}
+c_{j\sigma}^{\dagger}c_{i\sigma}\right),
\label{eq:app_hubbard_variants}
\end{align}
with $n_i=n_{i\uparrow}+n_{i\downarrow}$. The spinless fermion control family is
\begin{equation}
H_{tV}=
-t\sum_{\langle i,j\rangle}
\left(c_i^\dagger c_j+c_j^\dagger c_i\right)
+V\sum_{\langle i,j\rangle}n_i n_j
+\sum_i\epsilon_i n_i .
\label{eq:app_spinless_tv}
\end{equation}

The Hubbard--Holstein benchmark augments Eq.~(\ref{eq:app_hubbard}) with truncated phonons and density-coupled lattice displacement,
\begin{align}
H_{\hh}=&H_{\rm Hub}[\epsilon]
+\omega_0\sum_i b_i^{\dagger}b_i
+g\sum_i X_i(n_i-\bar n),
\qquad
X_i=b_i^\dagger+b_i .
\label{eq:app_hh}
\end{align}
For the half-filled lattice cases $\bar n=1$. Field offsets enter through the onsite coefficients $\epsilon_i$ in the Hubbard term.

The boson-only benchmark families use the same truncated local Fock operators. The number-conserving Bose--Hubbard chain is
\begin{align}
H_{\rm BH}={}&
-t\sum_{\langle i,j\rangle}
\left(b_i^\dagger b_j+b_j^\dagger b_i\right)
+\omega_0\sum_i n_i
\nonumber\\
&
+\frac{U_b}{2}\sum_i n_i(n_i-I)
+\sum_i\epsilon_i n_i ,
\label{eq:app_bose_hubbard}
\end{align}
while the harmonic/Kerr chain used for bosonic stress tests is
\begin{align}
H_{\rm HK}={}&
\omega_0\sum_i n_i
+\frac{U_K}{2}\sum_i n_i(n_i-I)
\nonumber\\
&
-J_X\sum_{\langle i,j\rangle}X_iX_j
+\sum_i\eta_iX_i .
\label{eq:app_harmonic_kerr}
\end{align}
The optional zero-point identity contribution is tracked in the benchmark specification and shifts only scalar energy offsets.

The spin-boson/generalized-Rabi control family is represented by a single two-level emitter coupled to one truncated boson mode,
\begin{align}
H_{\rm SB}={}&
-t_s\sigma_x+\Delta\sigma_z
+\omega_0 b^\dagger b
\nonumber\\
&
+u_xX_b\sigma_x
+g_zX_b\sigma_z,
\qquad X_b=b^\dagger+b .
\label{eq:app_spin_boson}
\end{align}
The restricted closed-shell molecular electronic tests use
\begin{equation}
H_{\rm mol}=E_{\rm nuc}I
+\sum_{pq}h_{pq}a_p^\dagger a_q
+\frac12\sum_{pqrs}(pr|qs)
a_p^\dagger a_q^\dagger a_s a_r .
\label{eq:app_molecular}
\end{equation}
Molecular vibronic and generalized electron--phonon instances add truncated modes to Eq.~(\ref{eq:app_molecular}) through
\begin{align}
H_{\rm vib}={}&H_{\rm mol}
+\sum_\mu\omega_\mu b_\mu^\dagger b_\mu
\nonumber\\
&
+\sum_{\mu,pq}
\left(\kappa^X_{\mu pq}X_\mu+\kappa^P_{\mu pq}P_\mu\right)
a_p^\dagger a_q
+H_b^{\rm nl}.
\label{eq:app_vibronic}
\end{align}

\section{SNAKE selector feature details}
\label{app:static_controller_details}

The selector-time variable $t$ denotes the SNAKE iteration index. Let $D_{\rm left}(t)$ be the remaining admission budget after depth, parameter, and stopping caps. With look-ahead index $q$,
\begin{equation}
\widehat N_{\rm rem}(t)=
\sum_{q=1}^{D_{\rm left}(t)}
S_t(q)Q_t(q).
\label{eq:expected_utility_left}
\end{equation}
Here $S_t(q)$ is the survival factor for a viable, budget-compatible admission opportunity $q$ selector steps ahead, and $Q_t(q)$ is the conditional usefulness factor for that opportunity. The normalized runway coordinate is
\begin{equation}
\begin{aligned}
R_t&=\frac{\widehat N_{\rm rem}(t)}{D_{\rm left}(t)+\eps}\in[0,1],\\
e_t&=R_t^a,\qquad
g_t=(1-R_t)^b,
\end{aligned}
\label{eq:runway_coordinate}
\end{equation}
with $a,b>0$. Thus $e_t$ weights early exploration and $g_t$ weights late insertion pressure.

For batch admission, let $\Gamma_{\rm batch}^{\rm proxy}(B;t)\in\{0,1\}$ denote the cheap proxy gate for joint insertability, support overlap, commutation, layout scheduling, and measurement reuse. After Phase 3 has produced $\Sset_3(t)$,
\begin{equation}
\mathfrak B_{\rm proxy}(t)=
\{B\subseteq\Sset_3(t):\Gamma_{\rm batch}^{\rm proxy}(B;t)=1\}.
\label{eq:batch_proxy_screen}
\end{equation}
Let $s_t(r)=S_3(r;t)$ and $r_{\max}\in\argmaxp_{r\in\Sset_3(t)}s_t(r)$. The near-degenerate shell and proxy-feasible batch family are
\begin{equation}
\begin{aligned}
\mathcal C_{\rm nd}(t)
&=\{r\in\Sset_3(t):s_t(r)\ge \eta_{\rm nd}s_t(r_{\max})\},\\
\mathfrak B_{\rm nd}(t)
&=\{B\subseteq\mathcal C_{\rm nd}(t):r_{\max}\in B,\ |B|\le B_{\max}(t)\},\\
\mathfrak B_{\rm geom}(t)
&=\mathfrak B_{\rm nd}(t)\cap\mathfrak B_{\rm proxy}(t).
\end{aligned}
\label{eq:near_degenerate_shell}
\end{equation}
For $B\in\mathfrak B_{\rm geom}(t)$, with proposed amplitudes $\alpha=(\alpha_r)_{r\in B}$, reduced Gram matrix $G_B$, reduced curvature $\widetilde H_B$, gain vector $g_{B,\rm lcb}$, and batch trust radius $\rho_B(t)$,
\begin{equation}
\begin{aligned}
\mathcal T_B(t)&=\{\alpha:\alpha^{\top}G_B(t)\alpha\le \rho_B(t)^2\},\\
\Delta E_B^{\TR}(t)={}&\max_{\alpha\in\mathcal T_B(t)}
\Bigl[g_{B,\rm lcb}^{\top}\alpha
-\frac12\alpha^{\top}(\widetilde H_B)^{(+)}\alpha\Bigr],\\
B_t^{\star}&\in\argmaxp_{B\in\mathfrak B_{\rm geom}(t)}\Delta E_B^{\TR}(t).
\end{aligned}
\label{eq:batch_gain}
\end{equation}

The generator-ablation guards include an amplitude-collapse witness. With $a_j(q)$ the stored normalized amplitude history and $a_{j,\rm early}^{\max}$, $a_{j,\rm late}^{\max}$ the maxima over the early and late windows,
\begin{equation}
\begin{aligned}
\Gamma_j^{\rm amp\downarrow}(t)=1
\iff{}&
a_{j,\rm early}^{\max}\ge a_{\rm hi},\\
&a_{j,\rm late}^{\max}\le a_{\rm lo},\\
&{\rm MAD}_{W_j^{\rm late}}(a_j)\le \sigma_{\rm lo}.
\end{aligned}
\label{eq:amp_collapse_witness}
\end{equation}
The cheap generator-ablation nomination set is
\begin{equation}
\begin{aligned}
\Dcal_{\rm abl}=\bigl\{j:\ &t-\iota_j\ge A_{\rm mat},\ c_j^{\rm cool}=0,\\
&|\theta_j|\le \theta_{\rm stale}\ \vee\ \Gamma_j^{\rm amp\downarrow}(t)=1,\\
&\Delta E_{\rm ablate}(j)\le \tau_E^{\rm abl},\ K_j^{\rm saved}\ge \tau_K\bigr\}.
\end{aligned}
\label{eq:static_ablation_candidates}
\end{equation}
Here $\iota_j$ is the admission step of coordinate $j$. The formal recoverability target is
\begin{equation}
\delta E_j^\star
:=
\min_{\xi\in\Theta_{\Oset_{-j}}}
E(\Oset_{-j},\xi)
-
E(\Oset,\theta),
\qquad
j\ \text{is redundant if}\ \delta E_j^\star\le\epsilon_{\rm pr}.
\label{eq:static_ablation_recoverability}
\end{equation}
The generator-ablation ladder approximates this condition with cheap nominations, cached ablation proxies, and then an explicit local refit; it is not a theorem that the proxy alone certifies deletion. A final ablation trial accepts deletion only after a local refit verifies small regression. The plateau stopping guard uses
\begin{equation}
\begin{aligned}
{\rm drop}(d)&=\Delta E_{\rm abs}(d-1)-\Delta E_{\rm abs}(d),\\
{\rm stop}\iff{}&d\ge d_{\min}\ \wedge\ {\rm drop}(d')<\tau_{\rm drop}\\
&\text{for } P\ \text{recent depths}.
\end{aligned}
\label{eq:plateau_stop}
\end{equation}

\section{Selector telemetry}
\label{app:controller_telemetry}

The useful-runway forecast in Eq.~(\ref{eq:expected_utility_left}) is computed from recent validated admission gains and current frontier flatness. For a valid past admission step $s$, define
\begin{equation}
\begin{aligned}
\Delta_{\rm adm}(s)
&=E(\theta^{(s-1)})-E(\theta^{(s)}),\\
\Delta_{\rm adm}^{+}(s)
&=\operatorname{asinh}\!\left(\frac{\Delta_{\rm adm}(s)}{\tau_\Delta}\right),
\end{aligned}
\label{eq:appendix_admission_gain}
\end{equation}
and let $\mathcal I_t^+$ be the recent index set used for EWMA/MAD statistics. The gain, decay, and scale summaries are
\begin{equation}
\begin{aligned}
m_t&=\operatorname{EWMA}_{s\in\mathcal I_t^+}(\Delta_{\rm adm}^{+}(s)),\\
\rho_t&=\operatorname{clip}\!\left(
\frac{\operatorname{EWMA}_{\rm recent}(\Delta_{\rm adm})}
{\operatorname{EWMA}_{\rm older}(\Delta_{\rm adm})+\eps},
\rho_{\min},1\right),\\
\gamma_t&=[-\log\rho_t]_+,\\
s_t&=\max\!\left\{s_{\min},
\operatorname{MAD}_{s\in\mathcal I_t^+}(\Delta_{\rm adm}^{+}(s))\right\}.
\end{aligned}
\label{eq:appendix_runway_stats}
\end{equation}
If $s_{{\rm raw},1}^{(1)}(t)\ge s_{{\rm raw},2}^{(1)}(t)$ are the two leading raw Phase-1 scores, the frontier-flatness statistic is
\begin{equation}
u_{\rm front}(t)=
\frac{s_{{\rm raw},2}^{(1)}(t)+\eps}
{s_{{\rm raw},1}^{(1)}(t)+\eps}.
\label{eq:appendix_frontier_flatness}
\end{equation}
Together with a stagnation statistic $u_{\rm stag}(t)$, this gives the survival hazard
\begin{equation}
\begin{split}
h_t(j)=\operatorname{sigm}\!\bigl(&
\beta_0+\beta_{\rm stag}u_{\rm stag}(t)\\
&+\beta_{\rm front}u_{\rm front}(t)+\eta_h(j-1)
\bigr),
\end{split}
\label{eq:appendix_survival_hazard}
\end{equation}
where $\operatorname{sigm}$ is the logistic sigmoid. The survival factor in Eq.~(\ref{eq:expected_utility_left}) is
\begin{equation}
S_t(k)=\prod_{j=1}^{k-1}\bigl(1-h_t(j)\bigr).
\label{eq:appendix_survival_usefulness}
\end{equation}
The usefulness probability is $Q_t(k)=\Phi(z_t(k))$, with
$z_t(k)=[m_t-\gamma_t(k-1)-u_t]/s_t$,
$u_t=\operatorname{asinh}(\tau_{{\rm use},t}/\tau_\Delta)$, and
$\Phi$ the standard normal cumulative distribution function.

The explicit shot schedule used behind the main-text prose is
\begin{equation}
\begin{aligned}
\Phi_k(t)&=g_t+(1-g_t)h_k(t),\qquad 0\le h_k(t)\le1,\\
N_{{\rm shot},k}^{\rm eff}(t)
&=\Gamma_k^{\rm live}(t)
\min\{N_{{\rm shot},k}^{\rm cap},N_{{\rm shot},k}^{\rm dec}\},\\
N_{{\rm shot},k}^{\rm dec}
&=\max\{N_{{\rm shot},k}^{\rm sched},N_{{\rm shot},k}^{\rm snr}\},\\
N_{{\rm shot},k}^{\rm sched}
&=\left\lceil N_{{\rm shot},k}^{\min}
+(N_{{\rm shot},k}^{\max}-N_{{\rm shot},k}^{\min})\Phi_k(t)\right\rceil,\\
N_{{\rm shot},k}^{\rm snr}
&=\left\lceil
\frac{\kappa_k^2\widehat\sigma_{t,k}^{\,2}}
{\max\{\widehat\delta_{t,k}^{\,2},\delta_{{\rm floor},k}^{\,2}\}}
\right\rceil .
\end{aligned}
\label{eq:appendix_shot_schedule}
\end{equation}
Here $h_k(t)$ is the phase-local uncertainty uplift, $\Gamma_k^{\rm live}$ is the live-stage mask, $\widehat\delta_{t,k}$ is the decision-scale signal, and $\widehat\sigma_{t,k}$ is its estimator spread. The matching shortlist and live-stage rules are
\begin{equation}
\begin{aligned}
\tau_k^{\rm short}(t)&=\tau_k^{\min}+(\tau_k^{\max}-\tau_k^{\min})g_t,\\
B_k^{\rm short}(t)&=
\left\lceil B_k^{\min}
+(B_k^{\max}-B_k^{\min})e_t\right\rceil,\\
\nu_{{\rm front},q}^{(\ell)}(t)
&=\frac{s_{q+1}^{(\ell)}(t)+\eps}{s_q^{(\ell)}(t)+\eps},\\
k_{\star}(t)&=1+\Gamma_2^{\rm live}(t)
+\Gamma_2^{\rm live}(t)\Gamma_3^{\rm live}(t).
\end{aligned}
\label{eq:appendix_shortlist_live_stage}
\end{equation}
Thus $R_t\downarrow$ increases insertion pressure, contracts the shortlist, and retires the expensive phases in the order Phase 3 then Phase 2.

\section{Evaluation details}
\label{app:evaluation_details}

The benchmark suite is stratified by Hamiltonian family, size, coupling, phonon cutoff, encoding, layout, and random seed, as summarized in Table~\ref{tab:benchmark_suite}. The same fixed SNAKE parameters $\phi_{\rm gen}^{\star}$ are used for the held-out rows.

\subsection{ADAPT descent trace}

The Hubbard--Holstein descent trace summarizes how each method purchases energy improvement along the admitted-generator axis. Several competitors are nearly flat over this axis, so the main resource comparison is read from the Pareto tables rather than from convergence speed alone.

\begin{center}
\centering
\IfFileExists{figures/hh_static_descent_tentative.png}{%
  \includegraphics[width=0.97\columnwidth]{figures/hh_static_descent_tentative.png}%
}{%
  \IfFileExists{hh_static_descent_tentative.png}{%
    \includegraphics[width=0.97\columnwidth]{hh_static_descent_tentative.png}%
  }{%
    \IfFileExists{../../plots/main_condensed_tentative_20260502/hh_static_descent_tentative.png}{%
      \includegraphics[width=0.97\columnwidth]{../../plots/main_condensed_tentative_20260502/hh_static_descent_tentative.png}%
    }{%
      \textsf{[Missing figure]}%
    }%
  }%
}
\end{center}
\noindent\emph{ADAPT descent figure.} The horizontal axis is SNAKE iteration or admitted-generator index, the vertical axis is ground-state energy error, and the overlays use the same Hubbard--Holstein point, cutoff, encoding, backend target, and seed.

\begin{table*}[t]
\small
\setlength{\tabcolsep}{4pt}
\caption{Benchmark suite for ADAPT ansatz-construction and resource-stress evaluation.}
\label{tab:benchmark_suite}
\begin{ruledtabular}
\begin{tabular*}{\textwidth}{@{\extracolsep{\fill}}p{0.17\textwidth}p{0.33\textwidth}ccp{0.23\textwidth}@{}}
Family & Representative parameters & Main rows & Encoding/cutoff panel & Held-out stress tests \\
\colrule
Hubbard / extended Hubbard & $L=\placeholder{2--6}$, $U/t=\placeholder{0,2,4,8}$, $V/t=\placeholder{0--2}$ & \placeholder{$N_1$} & \placeholder{$N_2$} & strong $U/t$, boundary changes \\
Hubbard--Holstein & $L=\placeholder{2--4}$, $g/t=\placeholder{0--1.5}$, $\omega_0/t=\placeholder{0.5--2}$, $\phmax=\placeholder{1--3}$ & \placeholder{$N_3$} & \placeholder{$N_4$} & larger cutoff and layout variation \\
Generalized electron--phonon & finite-range or bond-phonon coupling & \placeholder{$N_5$} & \placeholder{$N_6$} & SSH/Peierls-like coupling \\
Boson-only controls & Bose--Hubbard, harmonic/Kerr chain & \placeholder{$N_7$} & \placeholder{$N_8$} & high occupation cutoff \\
Molecular electronic & restricted closed-shell active spaces & \placeholder{$N_9$} & -- & stretched bonds / strong correlation \\
Molecular vibronic & H$_2$-like vibronic prototype, Anderson--Holstein analogue & \placeholder{$N_{10}$} & \placeholder{$N_{11}$} & displaced surfaces / multiple modes \\
\end{tabular*}
\end{ruledtabular}
\end{table*}

The reduced encoding panel varies the fermion map over JW, parity, and BK; the register layout over spin-blocked, site-interleaved fermions, and site-local fermion--boson blocks; and the boson encoding over binary, Gray, and unary registers. The size/cutoff pairs are $(L,\phmax)\in\{(2,1),(2,2),(2,3),(3,1),(3,2)\}$, with $U/t\in\{0,4,8\}$ and $g/t\in\{0.25,0.5,1.0\}$. Backend checks compare the canonical IBM Marrakesh/FakeMarrakesh heavy-hex target with simpler line, square, and all-to-all coupling proxies; the main tables use Marrakesh because the small-system ordering is unchanged within \textsf{[backend variation]}.

Ground-state accuracy is
\begin{equation}
\begin{aligned}
|\Delta E|&=|E_{\rm alg}-E_{\ex}|,\\
1-F&=1-|\braket{\psi_{\rm alg}}{\psi_{\ex}}|^2,\\
\Delta_O&=\expval{O}_{\rm alg}-\expval{O}_{\ex}.
\end{aligned}
\label{eq:appendix_static_metrics}
\end{equation}
for ground-state energy, state fidelity, and ground-state observable errors. Shot estimates use grouped Pauli measurement variances,
\begin{equation}
\Var(\widehat E)=\sum_{\alpha}\frac{h_{\alpha}^2\Var(P_{\alpha})}{N_{\alpha}},
\qquad
N_{\rm total}\propto |\Delta E|^{-2},
\label{eq:appendix_shot_estimate}
\end{equation}
or the grouped analogue used by the measurement allocator.

Noise robustness uses two lightweight checks. SNAKE supports analytic iid-Gaussian perturbations of measured energies, gradients, and metrics, with variance scaled by nominal shots and optional two-qubit-depth and measurement-group factors. The hardware-validation path also supports local fake-backend scheduled noise, readout mitigation, Pauli twirling, decoupling pulse sequences, and zero-noise extrapolation; those workflows serve as validation checks for the resource surface.

\section{Qubit encodings and register layout}
\label{app:encodings}

For the canonical Hubbard--Holstein benchmark, fermions are mapped by the Jordan--Wigner transformation. With blocked spin ordering $p(i,\uparrow)=i$ and $p(i,\downarrow)=L+i$,
\begin{equation}
\begin{aligned}
c_p^{\dagger}&=\frac12(X_p-\ii Y_p)\prod_{r=0}^{p-1}Z_r,\\
c_p&=\frac12(X_p+\ii Y_p)\prod_{r=0}^{p-1}Z_r,\\
n_p&=\frac{I-Z_p}{2} .
\end{aligned}
\label{eq:jw_map}
\end{equation}
Each local phonon mode is truncated to dimension $d=\phmax+1$ with occupation basis $\{\ket{n}_i\}_{n=0}^{d-1}$. In binary encoding,
\begin{equation}
\begin{aligned}
\ket{n}_i&\mapsto \ket{\operatorname{bin}(n)}_i,\\
b_i^{\dagger}&=\sum_{n=0}^{d-2}\sqrt{n+1}\,\ket{\operatorname{bin}(n+1)}_i\bra{\operatorname{bin}(n)}_i,\\
b_i&=\sum_{n=1}^{d-1}\sqrt{n}\,\ket{\operatorname{bin}(n-1)}_i\bra{\operatorname{bin}(n)}_i,\\
n_{b,i}&=\sum_{n=0}^{d-1}n\,\ket{\operatorname{bin}(n)}_i\bra{\operatorname{bin}(n)}_i,\\
x_i&=\sum_{n=0}^{d-2}\sqrt{n+1}\Bigl(\ket{\operatorname{bin}(n+1)}_i\bra{\operatorname{bin}(n)}_i\\
&\qquad\qquad\qquad\qquad+\ket{\operatorname{bin}(n)}_i\bra{\operatorname{bin}(n+1)}_i\Bigr),\\
P_i&=\ii\sum_{n=0}^{d-2}\sqrt{n+1}\Bigl(\ket{\operatorname{bin}(n+1)}_i\bra{\operatorname{bin}(n)}_i\\
&\qquad\qquad\qquad\qquad-\ket{\operatorname{bin}(n)}_i\bra{\operatorname{bin}(n+1)}_i\Bigr),
\end{aligned}
\label{eq:boson_binary_encoding}
\end{equation}
with the operator understood on the truncated binary code subspace. In unary encoding,
\begin{equation}
\begin{aligned}
\ket{n}_i&\mapsto \ket{e_n}_i\equiv\ket{0\cdots 010\cdots 0}_i,\\
b_i^{\dagger}&=\sum_{n=0}^{d-2}\sqrt{n+1}\,\ket{e_{n+1}}_i\bra{e_n}_i,\\
b_i&=\sum_{n=1}^{d-1}\sqrt{n}\,\ket{e_{n-1}}_i\bra{e_n}_i,\\
n_{b,i}&=\sum_{n=0}^{d-1}n\,\ket{e_n}_i\bra{e_n}_i,\\
x_i&=\sum_{n=0}^{d-2}\sqrt{n+1}\Bigl(\ket{e_{n+1}}_i\bra{e_n}_i\\
&\qquad\qquad\qquad\qquad+\ket{e_n}_i\bra{e_{n+1}}_i\Bigr),\\
P_i&=\ii\sum_{n=0}^{d-2}\sqrt{n+1}\Bigl(\ket{e_{n+1}}_i\bra{e_n}_i\\
&\qquad\qquad\qquad\qquad-\ket{e_n}_i\bra{e_{n+1}}_i\Bigr),
\end{aligned}
\label{eq:boson_unary_encoding}
\end{equation}
on the one-hot code subspace. The number of phonon qubits per site is
\begin{equation}
q_{\rm pb}=\begin{cases}
\max\{1,\lceil\log_2 d\rceil\}, & \text{binary},\\
d, & \text{unary},
\end{cases}
\label{eq:qpb}
\end{equation}
and the full Hubbard--Holstein register contains
\begin{equation}
N_q=2L+Lq_{\rm pb}.
\label{eq:nqubits}
\end{equation}
Gray-code and variational basis-state encodings are used only in the encoding/resource robustness study.

\section{Operator pools}
\label{app:pools}

The full operator pool is a deduplicated union
\begin{equation}
\Gset_{\rm full}=\Gset_{\rm UCCSD}\cup\Gset_{\rm HVA}\cup\Gset_{\rm term}^{\rm aug}\cup\Gset_{\rm PAOP}^{\rm full}\cup\Gset_{\rm PAOP,LF}^{\rm full}.
\label{eq:full_meta_pool}
\end{equation}
For the reduced Hubbard--Holstein tests, the lean family is
\begin{equation}
\Gset_{\rm lean}=\Gset_{\rm UCCSD}^{\rm lift}\cup\Gset_{\rm term}^{\rm quad}\cup\Gset_{\rm cloud,p}\cup\Gset_{\rm disp}\cup\Gset_{\rm hopdrag}\cup\Gset_{\rm dbl,p},
\label{eq:lean_pool}
\end{equation}
with a narrower $L=2$, $\phmax=1$ family
\begin{equation}
\begin{aligned}
\Gset_{{\rm lean},L=2}={}&\Gset_{\rm UCCSD,sing}^{\rm lift}\cup\Gset_{\rm term}^{\rm quad}\\
&\cup\Gset_{\rm cloud,p}\cup\Gset_{\rm hopdrag}\cup\Gset_{\rm dbl,p}^{\rm cross-site}.
\end{aligned}
\label{eq:l2_lean_pool}
\end{equation}
Representative symbolic generators are
\begin{align}
A_{ia,\sigma}^{\rm sing}&=\ii\left(c_{a\sigma}^{\dagger}c_{i\sigma}-c_{i\sigma}^{\dagger}c_{a\sigma}\right),
\nonumber\\
A_{ijab}^{\rm dbl}&=\ii\left(c_a^{\dagger}c_b^{\dagger}c_jc_i-\mathrm{h.c.}\right),
\nonumber\\
G_i^{({\rm cloud},p)}&=\tilde n_i\sum_{r\in\mathcal N(i)}\alpha_{ir}P_r,
\nonumber\\
G_{ij}^{({\rm hd})}&=T_{ij}^{(+)}(P_i-P_j),
\nonumber\\
G_i^{({\rm disp})}&\sim \tilde n_ix_i,
\nonumber\\
G_i^{({\rm dd},p)}&\sim D_iP_i,
\label{eq:pool_family_symbols}
\end{align}
where $x_i=b_i+b_i^{\dagger}$ and $P_i=\ii(b_i^{\dagger}-b_i)$.


\begin{thebibliography}{99}

\bibitem{Hubbard1963}
J. Hubbard,
``Electron correlations in narrow energy bands,''
\emph{Proceedings of the Royal Society of London A} \textbf{276}, 238 (1963).

\bibitem{Holstein1959a}
T. Holstein,
``Studies of polaron motion: Part I. The molecular-crystal model,''
\emph{Annals of Physics} \textbf{8}, 325 (1959).

\bibitem{Holstein1959b}
T. Holstein,
``Studies of polaron motion: Part II. The small polaron,''
\emph{Annals of Physics} \textbf{8}, 343 (1959).

\bibitem{Weisse2008}
A. Wei{\ss}e and H. Fehske,
``Exact diagonalization techniques,''
in \emph{Computational Many-Particle Physics},
edited by H. Fehske, R. Schneider, and A. Wei{\ss}e,
Lecture Notes in Physics Vol. 739 (Springer, Berlin, Heidelberg, 2008), pp. 529--544.

\bibitem{Troyer2005}
M. Troyer and U.-J. Wiese,
``Computational complexity and fundamental limitations to fermionic quantum Monte Carlo simulations,''
\emph{Physical Review Letters} \textbf{94}, 170201 (2005).

\bibitem{Schollwock2011}
U. Schollw\"ock,
``The density-matrix renormalization group in the age of matrix product states,''
\emph{Annals of Physics} \textbf{326}, 96 (2011).

\bibitem{Peruzzo2014}
A. Peruzzo, J. McClean, P. Shadbolt, M.-H. Yung, X.-Q. Zhou, P.~J. Love, A. Aspuru-Guzik, and J.~L. O'Brien,
``A variational eigenvalue solver on a photonic quantum processor,''
\emph{Nature Communications} \textbf{5}, 4213 (2014).

\bibitem{McClean2016}
J.~R. McClean, J. Romero, R. Babbush, and A. Aspuru-Guzik,
``The theory of variational hybrid quantum-classical algorithms,''
\emph{New Journal of Physics} \textbf{18}, 023023 (2016).

\bibitem{Kandala2017}
A. Kandala, A. Mezzacapo, K. Temme, M. Takita, M. Brink, J.~M. Chow, and J.~M. Gambetta,
``Hardware-efficient variational quantum eigensolver for small molecules and quantum magnets,''
\emph{Nature} \textbf{549}, 242 (2017).

\bibitem{McClean2018BarrenPlateaus}
J.~R. McClean, S. Boixo, V.~N. Smelyanskiy, R. Babbush, and H. Neven,
``Barren plateaus in quantum neural network training landscapes,''
\emph{Nature Communications} \textbf{9}, 4812 (2018).

\bibitem{Grimsley2019}
H.~R. Grimsley, S.~E. Economou, E. Barnes, and N.~J. Mayhall,
``An adaptive variational algorithm for exact molecular simulations on a quantum computer,''
\emph{Nature Communications} \textbf{10}, 3007 (2019).

\bibitem{Grimsley2023ADAPT}
H.~R. Grimsley, G.~S. Barron, E. Barnes, S.~E. Economou, and N.~J. Mayhall,
``Adaptive, problem-tailored variational quantum eigensolver mitigates rough parameter landscapes and barren plateaus,''
\emph{npj Quantum Information} \textbf{9}, 19 (2023).

\bibitem{Tang2021}
H.~L. Tang, V.~O. Shkolnikov, G.~S. Barron, H.~R. Grimsley, N.~J. Mayhall, E. Barnes, and S.~E. Economou,
``Qubit-ADAPT-VQE: An adaptive algorithm for constructing hardware-efficient ans\"atze on a quantum processor,''
\emph{PRX Quantum} \textbf{2}, 020310 (2021).

\bibitem{Yordanov2021}
Y.~S. Yordanov, V. Armaos, C.~H.~W. Barnes, and D.~R.~M. Arvidsson-Shukur,
``Qubit-excitation-based adaptive variational quantum eigensolver,''
\emph{Communications Physics} \textbf{4}, 228 (2021).

\bibitem{Feniou2023}
C. Feniou, M. Hassan, D. Traor\'e, E. Giner, Y. Maday, and J.-P. Piquemal,
``Overlap-ADAPT-VQE: Practical quantum chemistry on quantum computers via overlap-guided compact ans\"atze,''
\emph{Communications Physics} \textbf{6}, 192 (2023).

\bibitem{Anastasiou2024}
P.~G. Anastasiou, Y. Chen, N.~J. Mayhall, E. Barnes, and S.~E. Economou,
``TETRIS-ADAPT-VQE: An adaptive algorithm that yields shallower, denser circuit ans\"atze,''
\emph{Physical Review Research} \textbf{6}, 013254 (2024).

\bibitem{Ramoa2025}
M. Ram\^oa, P.~G. Anastasiou, L.~P. Santos, N.~J. Mayhall, E. Barnes, and S.~E. Economou,
``Reducing the resources required by ADAPT-VQE using coupled exchange operators and improved subroutines,''
\emph{npj Quantum Information} \textbf{11}, 86 (2025).

\bibitem{Ramoa2025HessianRecycling}
M. Ramoa et al.,
``Reducing measurement costs by recycling the Hessian in adaptive variational quantum algorithms,''
\emph{Quantum Science and Technology} \textbf{10}, 015031 (2025).

\bibitem{Nocedal2006TrustRegion}
J. Nocedal and S. J. Wright,
``Trust-Region Methods,''
in \emph{Numerical Optimization}, Springer (New York, NY, 1994).

\bibitem{Ikhtiarudin2025ShotEfficient}
A. Ikhtiarudin, G. K. Sunnardianto, F. Fathurrahman, M. K. Agusta, and H. Dipojono,
``Shot-Efficient ADAPT-VQE via Reused Pauli Measurements and Variance-Based Shot Allocation,''
\emph{arXiv:2507.16879} (2025).

\bibitem{Sohail2026GeoADAPT}
M.~A. Sohail and T. Koike-Akino,
``Geo-ADAPT-VQE: Quantum Information Metric-Aware Circuit Optimization for Quantum Chemistry,''
\emph{arXiv:2603.10325} (2026).

\bibitem{Stadelmann2025GradientTroughs}
T. Stadelmann and others,
``Strategies for overcoming gradient troughs in the ADAPT-VQE algorithm,''
\emph{arXiv:2512.25004} (2025).

\bibitem{ProvostVallee1980}
J.-P. Provost and G. Vall\'ee,
``Riemannian structure on manifolds of quantum states,''
\emph{Communications in Mathematical Physics} \textbf{76}, 289 (1980).

\bibitem{Stokes2020QNG}
J. Stokes, J. Izaac, N. Killoran, and G. Carleo,
``Quantum natural gradient,''
\emph{Quantum} \textbf{4}, 269 (2020).

\bibitem{Verteletskyi2020CliqueCover}
V. Verteletskyi, T.-C. Yen, and A.~F. Izmaylov,
``Measurement optimization in the variational quantum eigensolver using a minimum clique cover,''
\emph{Journal of Chemical Physics} \textbf{152}, 124114 (2020).

\bibitem{Yen2020CompatibleMeasurements}
T.-C. Yen, V. Verteletskyi, and A.~F. Izmaylov,
``Measuring all compatible operators in one series of single-qubit measurements using unitary transformations,''
\emph{Journal of Chemical Theory and Computation} \textbf{16}, 2400 (2020).


\bibitem{Macridin2018}
A. Macridin, P. Spentzouris, J. Amundson, and R. Harnik,
``Electron-phonon systems on a universal quantum computer,''
\emph{Physical Review Letters} \textbf{121}, 110504 (2018).

\bibitem{Sawaya2020}
N.~P.~D. Sawaya, T. Menke, T.~H. Kyaw, S. Johri, A. Aspuru-Guzik, and G.~G. Guerreschi,
``Resource-efficient digital quantum simulation of $d$-level systems for photonic, vibrational, and spin-$s$ Hamiltonians,''
\emph{npj Quantum Information} \textbf{6}, 49 (2020).

\bibitem{Li2023}
W. Li, J. Ren, S. Huai, T. Cai, Z. Shuai, and S. Zhang,
``Efficient quantum simulation of electron-phonon systems by variational basis state encoder,''
\emph{Physical Review Research} \textbf{5}, 023046 (2023).

\bibitem{Denner2023}
M.~M. Denner, H. Yan, N. A. Sinitsyn, and others,
``A hybrid quantum-classical method for electron-phonon systems,''
\emph{Communications Physics} \textbf{6}, 317 (2023).

\bibitem{Dutta2024BosonicChemistry}
R. Dutta et al.,
``Simulating chemistry on bosonic quantum devices,''
\emph{Journal of Chemical Theory and Computation} \textbf{20}, 6426--6441 (2024).

\bibitem{Peng2025BosonHamiltonians}
B. Peng, Y. Su, D. Claudino, K. Kowalski, G. H. Low, and M. Roetteler,
``Quantum simulation of boson-related Hamiltonians: Techniques, effective Hamiltonian construction, and error analysis,''
\emph{Quantum Science and Technology} \textbf{10}, 023002 (2025).

\bibitem{Vaquero2025PrunedADAPT}
N. Vaquero-Sabater, A. Carreras, and D. Casanova,
``Pruned-ADAPT-VQE: Compacting Molecular Ans\"atze by Removing Irrelevant Operators,''
\emph{Journal of Chemical Theory and Computation} \textbf{21}, 8720--8728 (2025).

\end{thebibliography}

% ---------- pasteable proof block ----------
\section{Convergence of the Single-Record Reduced-Geometry Selector}

We prove a convergence statement for the idealized infinite-run,
single-record form of the Phase-3 selector.  The proof uses the Phase-3
objects already defined above: the candidate record set
\(\mathcal R_3(T)\), reduced-window quantities \(F_r^{\mathrm{red}}\),
\(F_{r,\mathrm{safe}}^{\mathrm{red}}\), \(q_r^{\mathrm{red}}\), reduced
novelty \(\mathcal N_3(r;T)\), trust-region gain
\(\Delta E_{\mathrm{TR}}(r)\), burden \(K_3(r;T)\), and score
\(S_3(r;T)\).  Here \(T=0,1,2,\ldots\) denotes the selector step in this
convergence argument.

At step \(T\), let the current scaffold and refitted parameters be
\((\mathcal O_T,\theta_T)\), and define the actual variational energy
\[
E_T:=E(\mathcal O_T,\theta_T).
\]
For nonterminal steps, the single-record Phase-3 admission rule is
\[
r_T^\star=(m_T^\star,p_T^\star)
\in
\arg\max_{r\in\mathcal R_3(T)}S_3(r;T),
\qquad
S_3(r;T)
=
\frac{\Delta E_{\mathrm{TR}}(r)\mathcal N_3(r;T)}
{1+K_3(r;T)}.
\]
The next scaffold is
\[
\mathcal O_{T+1}:=\mathcal O_T\oplus_{p_T^\star}m_T^\star,
\]
and \(\theta_{T+1}\in\Theta_{\mathcal O_{T+1}}\) denotes the parameter
vector returned by the post-admission refit.

For the physical-geometric stationarity residual, define
\[
\Xi_T
:=
\sup_{r\in\mathcal R_3(T)}
\Delta E_{\mathrm{TR}}(r)\mathcal N_3(r;T).
\]
Thus \(\Xi_T=0\) means that every surviving Phase-3 record has zero
reduced trust-region gain, zero reduced novelty, or both.

\begin{schurlemma}[Reduced-window Schur identity]
For a record \(r=(m,p)\in\mathcal R_3(T)\), set
\[
\eta_r:=\mathbf R_r^{-1}b_r.
\]
Then the reduced norm appearing in Phase 3 satisfies
\[
F_r^{\mathrm{red}}
=
F_{\mathrm{raw}}(r)
-2q_r^\top\eta_r
+\eta_r^\top Q_r\eta_r.
\]
Equivalently, up to the already included \(\varepsilon_F\) in
\(F_{\mathrm{raw}}\),
\[
F_r^{\mathrm{red}}
=
\left\|
\tau_r-\sum_{j\in W_r}(\eta_r)_j\tau_j
\right\|^2+\varepsilon_F.
\]
Moreover,
\[
F_{r,\mathrm{safe}}^{\mathrm{red}}\mathcal N_3(r;T)
=
\left[
F_{r,\mathrm{safe}}^{\mathrm{red}}
-
(q_r^{\mathrm{red}})^\top
(Q_r+\varepsilon_{\mathrm{nov}}I)^{-1}
q_r^{\mathrm{red}}
\right]_{+}.
\]
Hence \(\mathcal N_3(r;T)\) is exactly the normalized reduced Schur
residual after energy-relaxed reduction against the inherited refit
window.
\end{schurlemma}

\begin{proof}
The first identity is only the definition of \(F_r^{\mathrm{red}}\) with
\(\eta_r=\mathbf R_r^{-1}b_r\).  For the norm interpretation, use the
Gram definitions
\[
(Q_r)_{jk}=\Re\langle\tau_j,\tau_k\rangle,
\qquad
(q_r)_j=\Re\langle\tau_j,\tau_r\rangle.
\]
Then
\[
\begin{aligned}
\left\|
\tau_r-\sum_{j\in W_r}(\eta_r)_j\tau_j
\right\|^2+\varepsilon_F
&=
\langle\tau_r,\tau_r\rangle+\varepsilon_F
-2q_r^\top\eta_r+\eta_r^\top Q_r\eta_r\\
&=
F_{\mathrm{raw}}(r)-2q_r^\top\eta_r+\eta_r^\top Q_r\eta_r\\
&=
F_r^{\mathrm{red}}.
\end{aligned}
\]
Finally, by definition,
\[
z_r
=
\frac{
(q_r^{\mathrm{red}})^\top
(Q_r+\varepsilon_{\mathrm{nov}}I)^{-1}
q_r^{\mathrm{red}}
}{
F_{r,\mathrm{safe}}^{\mathrm{red}}
},
\qquad
\mathcal N_3(r;T)=\operatorname{clip}_{[0,1]}(1-z_r).
\]
Since \(Q_r+\varepsilon_{\mathrm{nov}}I\succ0\), the numerator is
nonnegative, hence \(z_r\ge0\).  Therefore
\[
F_{r,\mathrm{safe}}^{\mathrm{red}}\mathcal N_3(r;T)
=
F_{r,\mathrm{safe}}^{\mathrm{red}}[1-z_r]_{+}
=
\left[
F_{r,\mathrm{safe}}^{\mathrm{red}}
-
(q_r^{\mathrm{red}})^\top
(Q_r+\varepsilon_{\mathrm{nov}}I)^{-1}
q_r^{\mathrm{red}}
\right]_{+}.
\]
\end{proof}

\begin{schurassumption}[Finite burden or disabled burden]
There exists \(C_K<\infty\) such that, for every nonterminal selector
step \(T\) and every \(r\in\mathcal R_3(T)\),
\[
1\le 1+K_3(r;T)\le C_K.
\]
If cost is disabled in the convergence analysis, take \(K_3(r;T)\equiv0\)
and hence \(C_K=1\).
\end{schurassumption}

\begin{schurassumption}[Actual-energy refit descent]
There exists \(c_{\mathrm{dec}}>0\) and a nonnegative summable sequence
\(\{\varepsilon_T^{\mathrm{refit}}\}_{T\ge0}\) such that
\[
\sum_{T=0}^{\infty}\varepsilon_T^{\mathrm{refit}}<\infty
\]
and every nonterminal single-record admission satisfies
\[
E_{T+1}
\le
E_T
-
 c_{\mathrm{dec}}\,
\Delta E_{\mathrm{TR}}(r_T^\star)\mathcal N_3(r_T^\star;T)
+
\varepsilon_T^{\mathrm{refit}}.
\]
This assumption is automatic if the implementation uses an accept/reject
rule with this sufficient decrease.  Without such a rule, it is an
assumption on the post-admission refit.
\end{schurassumption}

\begin{schurassumption}[Lower energy bound]
There exists \(E_{\inf}>-\infty\) such that
\[
E_T\ge E_{\inf}
\qquad
\text{for all }T.
\]
\end{schurassumption}

\begin{schurlemma}[Selected gain controls the Phase-3 residual]
Under finite burden,
\[
\Delta E_{\mathrm{TR}}(r_T^\star)\mathcal N_3(r_T^\star;T)
\ge
\frac{1}{C_K}\,\Xi_T.
\]
\end{schurlemma}

\begin{proof}
For compactness, write
\[
A_T(r):=\Delta E_{\mathrm{TR}}(r)\mathcal N_3(r;T).
\]
Since \(r_T^\star\) maximizes \(S_3(\cdot;T)\),
\[
\frac{A_T(r_T^\star)}{1+K_3(r_T^\star;T)}
\ge
\frac{A_T(r)}{1+K_3(r;T)}
\qquad
\forall r\in\mathcal R_3(T).
\]
Using \(1+K_3(r;T)\le C_K\) gives
\[
\frac{A_T(r_T^\star)}{1+K_3(r_T^\star;T)}
\ge
\frac{A_T(r)}{C_K}.
\]
Since \(1+K_3(r_T^\star;T)\ge1\),
\[
A_T(r_T^\star)
\ge
\frac{A_T(r_T^\star)}{1+K_3(r_T^\star;T)}
\ge
\frac{A_T(r)}{C_K}.
\]
Taking the supremum over \(r\in\mathcal R_3(T)\) yields the claim.
\end{proof}

\begin{schurtheorem}[Actual-energy convergence and Schur stationarity]
Assume finite burden, actual-energy refit descent, and a lower energy
bound.  Then the actual energy sequence \(E_T\) converges to a finite
limit \(E_\infty\ge E_{\inf}\), and
\[
\sum_{T=0}^{\infty}\Xi_T<\infty,
\qquad
\Xi_T\to0.
\]
Consequently the infinite-run single-record Phase-3 selector is
Schur-stationary in the following sense:
\[
\sup_{r\in\mathcal R_3(T)}
\Delta E_{\mathrm{TR}}(r)\mathcal N_3(r;T)
\longrightarrow0.
\]
Thus, asymptotically, no surviving Phase-3 candidate is simultaneously
predicted to be energy-useful and reduced-Schur-novel relative to the
current inherited refit window.
\end{schurtheorem}

\begin{proof}
Let
\[
\Omega_T:=\sum_{s=T}^{\infty}\varepsilon_s^{\mathrm{refit}}.
\]
Since the refit errors are summable, \(\Omega_T<\infty\) and
\(\Omega_T\to0\).  From actual-energy descent and
\(\Omega_T=\varepsilon_T^{\mathrm{refit}}+\Omega_{T+1}\),
\[
\begin{aligned}
E_{T+1}+\Omega_{T+1}
&\le
E_T
-
c_{\mathrm{dec}}\,
\Delta E_{\mathrm{TR}}(r_T^\star)\mathcal N_3(r_T^\star;T)
+\varepsilon_T^{\mathrm{refit}}
+\Omega_{T+1}\\
&=
E_T+\Omega_T
-
c_{\mathrm{dec}}\,
\Delta E_{\mathrm{TR}}(r_T^\star)\mathcal N_3(r_T^\star;T).
\end{aligned}
\]
Hence \(E_T+\Omega_T\) is monotone nonincreasing.  It is bounded below
by \(E_{\inf}\), so it converges to a finite limit.  Since
\(\Omega_T\to0\), the sequence \(E_T\) also converges to a finite limit,
denoted \(E_\infty\ge E_{\inf}\).

Summing the same inequality from \(T=0\) to \(N\) gives
\[
c_{\mathrm{dec}}
\sum_{T=0}^{N}
\Delta E_{\mathrm{TR}}(r_T^\star)\mathcal N_3(r_T^\star;T)
\le
E_0+\Omega_0-E_{N+1}-\Omega_{N+1}
\le
E_0+\Omega_0-E_{\inf}.
\]
Therefore
\[
\sum_{T=0}^{\infty}
\Delta E_{\mathrm{TR}}(r_T^\star)\mathcal N_3(r_T^\star;T)
<\infty.
\]
Using the selected-gain lower bound,
\[
\frac{1}{C_K}\sum_{T=0}^{\infty}\Xi_T
\le
\sum_{T=0}^{\infty}
\Delta E_{\mathrm{TR}}(r_T^\star)\mathcal N_3(r_T^\star;T)
<\infty,
\]
so \(\sum_T\Xi_T<\infty\).  Since each \(\Xi_T\ge0\), summability implies
\(\Xi_T\to0\).
\end{proof}

\begin{schurassumption}[Schur-QGI / reduced-geometry gradient domination]
Let \(E_\star\le E_T\) denote the target energy infimum over the
reachable scaffold class.  There exists \(\mu_{\mathrm{SQGI}}>0\) such
that every nonterminal selector step satisfies
\[
\Xi_T
\ge
\mu_{\mathrm{SQGI}}(E_T-E_\star).
\]
Assume additionally that
\[
\nu:=
\frac{c_{\mathrm{dec}}\mu_{\mathrm{SQGI}}}{C_K}
\in(0,1].
\]
\end{schurassumption}

\begin{schurlemma}[Perturbed linear recursion]
Let \(d_T\ge0\) and suppose
\[
d_{T+1}\le(1-\nu)d_T+\varepsilon_T
\]
for some \(\nu\in(0,1]\) and some nonnegative summable sequence
\(\{\varepsilon_T\}_{T\ge0}\).  Then
\[
d_T
\le
(1-\nu)^T d_0
\sum_{s=0}^{T-1}(1-\nu)^{T-1-s}\varepsilon_s,
\]
and the convolution term tends to zero.  Hence \(d_T\to0\).
\end{schurlemma}

\begin{proof}
Iterating the recursion gives
\[
d_T
\le
(1-\nu)^T d_0
\sum_{s=0}^{T-1}(1-\nu)^{T-1-s}\varepsilon_s.
\]
The first term tends to zero.  For the second term, fix
\(\epsilon>0\).  Since \(\sum_s\varepsilon_s<\infty\), choose \(T_0\)
such that \(\sum_{s=T_0}^{\infty}\varepsilon_s<\epsilon\).  Split the
convolution into
\[
\sum_{s=0}^{T-1}(1-\nu)^{T-1-s}\varepsilon_s
=
\sum_{s=0}^{T_0-1}(1-\nu)^{T-1-s}\varepsilon_s
+\sum_{s=T_0}^{T-1}(1-\nu)^{T-1-s}\varepsilon_s.
\]
The first sum is finite and tends to zero as \(T\to\infty\); the second
is bounded by \(\sum_{s=T_0}^{\infty}\varepsilon_s<\epsilon\).  Since
\(\epsilon\) was arbitrary, the convolution tends to zero.
\end{proof}

\begin{schurtheorem}[Conditional global convergence]
Assume the hypotheses of the actual-energy convergence theorem and the
Schur-QGI condition.  Then
\[
E_T\to E_\star.
\]
More precisely,
\[
E_T-E_\star
\le
(1-\nu)^T(E_0-E_\star)
+
e_T,
\]
where
\[
e_T:=
\sum_{s=0}^{T-1}
(1-\nu)^{T-1-s}\varepsilon_s^{\mathrm{refit}},
\qquad
e_T\to0.
\]
In the exact-refit case
\(\varepsilon_T^{\mathrm{refit}}\equiv0\), this reduces to the geometric
bound
\[
E_T-E_\star
\le
(1-\nu)^T(E_0-E_\star).
\]
\end{schurtheorem}

\begin{proof}
By actual-energy descent and the selected-gain lower bound,
\[
\begin{aligned}
E_{T+1}
&\le
E_T
-
\frac{c_{\mathrm{dec}}}{C_K}\Xi_T
+
\varepsilon_T^{\mathrm{refit}}.
\end{aligned}
\]
Using Schur-QGI,
\[
E_{T+1}
\le
E_T
-
\frac{c_{\mathrm{dec}}\mu_{\mathrm{SQGI}}}{C_K}
(E_T-E_\star)
+
\varepsilon_T^{\mathrm{refit}}.
\]
With \(d_T:=E_T-E_\star\) and
\(\nu=c_{\mathrm{dec}}\mu_{\mathrm{SQGI}}/C_K\), this is
\[
d_{T+1}\le(1-\nu)d_T+\varepsilon_T^{\mathrm{refit}}.
\]
The perturbed linear recursion lemma gives the stated bound and
\(d_T\to0\).
\end{proof}

\begin{schurremark}[Interpretation]
The local convergence theorem does not require diagonal dominance of a
pool metric and does not require inversion of a raw, possibly singular
pool covariance matrix.  The role of diagonal dominance is replaced by
the computable reduced Schur residual
\[
F_{r,\mathrm{safe}}^{\mathrm{red}}\mathcal N_3(r;T)
=
\left[
F_{r,\mathrm{safe}}^{\mathrm{red}}
-
(q_r^{\mathrm{red}})^\top
(Q_r+\varepsilon_{\mathrm{nov}}I)^{-1}
q_r^{\mathrm{red}}
\right]_{+}.
\]
The price of claiming global convergence is the Schur-QGI condition:
the best surviving reduced-Schur gain must dominate the current
suboptimality \(E_T-E_\star\).  Without that condition, the theorem
still gives actual-energy convergence and vanishing Phase-3
Schur-stationarity residual \(\Xi_T\to0\), but not necessarily global
optimality.
\end{schurremark}
\iffalse
\section{Convergence of the Single-Record Reduced-Geometry Selector}

We prove a convergence statement for the idealized infinite-run,
single-record form of the Phase-3 selector.  The proof uses the Phase-3
objects already defined above: the candidate record set
\(\mathcal R_3(T)\), reduced-window quantities
$F_r^{\mathrm{red}}$, $F_{r,\mathrm{safe}}^{\mathrm{red}}$,
$q_r^{\mathrm{red}}$, reduced novelty \(\mathcal N_3(r;T)\), trust-region
gain \(\Delta E_{\mathrm{TR}}(r)\), burden \(K_3(r;T)\), and score
\(S_3(r;T)\).  Here \(T=0,1,2,\ldots\) denotes the selector step in
this convergence argument.

At step \(T\), let the current scaffold and refitted parameters be
\((\mathcal O_T,\theta_T)\), and define the actual variational energy
\[
E_T:=E(\mathcal O_T,\theta_T).
\]
For nonterminal steps, the single-record Phase-3 admission rule is
\[
r_T^\star=(m_T^\star,p_T^\star)
\in
\operatorname*{arg\,max}_{r\in\mathcal R_3(T)}S_3(r;T),
\qquad
S_3(r;T)
=
\frac{\Delta E_{\mathrm{TR}}(r)\mathcal N_3(r;T)}
{1+K_3(r;T)}.
\]
The next scaffold is
\[
\mathcal O_{T+1}:=\mathcal O_T\oplus_{p_T^\star}m_T^\star,
\]
and \(\theta_{T+1}\in\Theta_{\mathcal O_{T+1}}\) denotes the parameter
vector returned by the post-admission refit.

For the physical-geometric stationarity residual, define
\[
\Xi_T
:=
\sup_{r\in\mathcal R_3(T)}
\Delta E_{\mathrm{TR}}(r)\mathcal N_3(r;T).
\]
Thus \(\Xi_T=0\) means that every surviving Phase-3 record has zero
reduced trust-region gain, zero reduced novelty, or both.

\begin{schurlemma}[Reduced-window Schur identity]
For a record \(r=(m,p)\in\mathcal R_3(T)\), set
\[
\eta_r:=\mathbf R_r^{-1}b_r.
\]
Then the reduced norm appearing in Phase 3 satisfies
\[
F_r^{\mathrm{red}}
=
F_{\mathrm{raw}}(r)
-2q_r^\top\eta_r
+\eta_r^\top Q_r\eta_r.
\]
Equivalently, up to the already included \(\varepsilon_F\) in
\(F_{\mathrm{raw}}\),
\[
F_r^{\mathrm{red}}
=
\left\|\n	au_r-\sum_{j\in W_r}(\eta_r)_j\tau_j\right\|^2+\varepsilon_F.
\]
Moreover,
\[
F_{r,\mathrm{safe}}^{\mathrm{red}}\mathcal N_3(r;T)
=
\left[
F_{r,\mathrm{safe}}^{\mathrm{red}}
-
(q_r^{\mathrm{red}})^\top
(Q_r+\varepsilon_{\mathrm{nov}}I)^{-1}
q_r^{\mathrm{red}}
\right]_{+}.
\]
Hence \(\mathcal N_3(r;T)\) is exactly the normalized reduced Schur
residual after energy-relaxed reduction against the inherited refit
window.
\end{schurlemma}

\begin{proof}
The first identity is only the definition of \(F_r^{\mathrm{red}}\) with
\(\eta_r=\mathbf R_r^{-1}b_r\).  For the norm interpretation, use the
Gram definitions
\[
(Q_r)_{jk}=\Re\langle\tau_j,\tau_k\rangle,
\qquad
(q_r)_j=\Re\langle\tau_j,\tau_r\rangle.
\]
Then
\[
\begin{aligned}
\left\|
\tau_r-\sum_{j\in W_r}(\eta_r)_j\tau_j
\right\|^2+\varepsilon_F
&=
\langle\tau_r,\tau_r\rangle+\varepsilon_F
-2q_r^\top\eta_r+\eta_r^\top Q_r\eta_r\\
&=
F_{\mathrm{raw}}(r)-2q_r^\top\eta_r+\eta_r^\top Q_r\eta_r\\
&=
F_r^{\mathrm{red}}.
\end{aligned}
\]
Finally, by definition,
\[
z_r
=
\frac{
(q_r^{\mathrm{red}})^\top
(Q_r+\varepsilon_{\mathrm{nov}}I)^{-1}
q_r^{\mathrm{red}}
}{
F_{r,\mathrm{safe}}^{\mathrm{red}}
},
\qquad
\mathcal N_3(r;T)=\operatorname{clip}_{[0,1]}(1-z_r).
\]
Since \(Q_r+\varepsilon_{\mathrm{nov}}I\succ0\), the numerator is
nonnegative, hence \(z_r\ge0\).  Therefore
\[
F_{r,\mathrm{safe}}^{\mathrm{red}}\mathcal N_3(r;T)
=
F_{r,\mathrm{safe}}^{\mathrm{red}}[1-z_r]_{+}
=
\left[
F_{r,\mathrm{safe}}^{\mathrm{red}}
-
(q_r^{\mathrm{red}})^\top
(Q_r+\varepsilon_{\mathrm{nov}}I)^{-1}
q_r^{\mathrm{red}}
\right]_{+}.
\]
\end{proof}

\begin{schurassumption}[Finite burden or disabled burden]
There exists \(C_K<\infty\) such that, for every nonterminal selector
step \(T\) and every \(r\in\mathcal R_3(T)\),
\[
1\le 1+K_3(r;T)\le C_K.
\]
If cost is disabled in the convergence analysis, take \(K_3(r;T)\equiv0\)
and hence \(C_K=1\).
\end{schurassumption}

\begin{schurassumption}[Actual-energy refit descent]
There exists \(c_{\mathrm{dec}}>0\) and a nonnegative summable sequence
\(\{\varepsilon_T^{\mathrm{refit}}\}_{T\ge0}\) such that
\[
\sum_{T=0}^{\infty}\varepsilon_T^{\mathrm{refit}}<\infty
\]
and every nonterminal single-record admission satisfies
\[
E_{T+1}
\le
E_T
-
c_{
\mathrm{dec}}\,
\Delta E_{\mathrm{TR}}(r_T^\star)\mathcal N_3(r_T^\star;T)
+
\varepsilon_T^{\mathrm{refit}}.
\]
This assumption is automatic if the implementation uses an accept/reject
rule with this sufficient decrease.  Without such a rule, it is an
assumption on the post-admission refit.
\end{schurassumption}

\begin{schurassumption}[Lower energy bound]
There exists \(E_{\inf}>-\infty\) such that
\[
E_T\ge E_{\inf}
\qquad
\text{for all }T.
\]
\end{schurassumption}

\begin{schurlemma}[Selected gain controls the Phase-3 residual]
Under finite burden,
\[
\Delta E_{\mathrm{TR}}(r_T^\star)\mathcal N_3(r_T^\star;T)
\ge
\frac{1}{C_K}\,\Xi_T.
\]
\end{schurlemma}

\begin{proof}
For compactness, write
\[
A_T(r):=\Delta E_{\mathrm{TR}}(r)\mathcal N_3(r;T).
\]
Since \(r_T^\star\) maximizes \(S_3(\cdot;T)\),
\[
\frac{A_T(r_T^\star)}{1+K_3(r_T^\star;T)}
\ge
\frac{A_T(r)}{1+K_3(r;T)}
\qquad
\forall r\in\mathcal R_3(T).
\]
Using \(1+K_3(r;T)\le C_K\) gives
\[
\frac{A_T(r_T^\star)}{1+K_3(r_T^\star;T)}
\ge
\frac{A_T(r)}{C_K}.
\]
Since \(1+K_3(r_T^\star;T)\ge1\),
\[
A_T(r_T^\star)
\ge
\frac{A_T(r_T^\star)}{1+K_3(r_T^\star;T)}
\ge
\frac{A_T(r)}{C_K}.
\]
Taking the supremum over \(r\in\mathcal R_3(T)\) yields the claim.
\end{proof}

\begin{schurtheorem}[Actual-energy convergence and Schur stationarity]
Assume finite burden, actual-energy refit descent, and a lower energy
bound.  Then the actual energy sequence \(E_T\) converges to a finite
limit \(E_\infty\ge E_{\inf}\), and
\[
\sum_{T=0}^{\infty}\Xi_T<\infty,
\qquad
\Xi_T\to0.
\]
Consequently the infinite-run single-record Phase-3 selector is
Schur-stationary in the following sense:
\[
\sup_{r\in\mathcal R_3(T)}
\Delta E_{\mathrm{TR}}(r)\mathcal N_3(r;T)
\longrightarrow0.
\]
Thus, asymptotically, no surviving Phase-3 candidate is simultaneously
predicted to be energy-useful and reduced-Schur-novel relative to the
current inherited refit window.
\end{schurtheorem}

\begin{proof}
Let
\[
\Omega_T:=\sum_{s=T}^{\infty}\varepsilon_s^{\mathrm{refit}}.
\]
Since the refit errors are summable, \(\Omega_T<\infty\) and
\(\Omega_T\to0\).  From actual-energy descent and
\(\Omega_T=\varepsilon_T^{\mathrm{refit}}+\Omega_{T+1}\),
\[
\begin{aligned}
E_{T+1}+\Omega_{T+1}
&\le
E_T
-
\c_{
\mathrm{dec}}\,
\Delta E_{\mathrm{TR}}(r_T^\star)\mathcal N_3(r_T^\star;T)
+\varepsilon_T^{\mathrm{refit}}
+\Omega_{T+1}\\
&=
E_T+\Omega_T
-
\c_{
\mathrm{dec}}\,
\Delta E_{\mathrm{TR}}(r_T^\star)\mathcal N_3(r_T^\star;T).
\end{aligned}
\]
Hence \(E_T+\Omega_T\) is monotone nonincreasing.  It is bounded below
by \(E_{\inf}\), so it converges to a finite limit.  Since
\(\Omega_T\to0\), the sequence \(E_T\) also converges to a finite limit,
denoted \(E_\infty\ge E_{\inf}\).

Summing the same inequality from \(T=0\) to \(N\) gives
\[
\c_{
\mathrm{dec}}
\sum_{T=0}^{N}
\Delta E_{\mathrm{TR}}(r_T^\star)\mathcal N_3(r_T^\star;T)
\le
E_0+\Omega_0-E_{N+1}-\Omega_{N+1}
\le
E_0+\Omega_0-E_{\inf}.
\]
Therefore
\[
\sum_{T=0}^{\infty}
\Delta E_{\mathrm{TR}}(r_T^\star)\mathcal N_3(r_T^\star;T)
<\infty.
\]
Using the selected-gain lower bound,
\[
\frac{1}{C_K}\sum_{T=0}^{\infty}\Xi_T
\le
\sum_{T=0}^{\infty}
\Delta E_{\mathrm{TR}}(r_T^\star)\mathcal N_3(r_T^\star;T)
<\infty,
\]
so \(\sum_T\Xi_T<\infty\).  Since each \(\Xi_T\ge0\), summability implies
\(\Xi_T\to0\).
\end{proof}

\begin{schurassumption}[Schur-QGI / reduced-geometry gradient domination]
Let \(E_\star\le E_T\) denote the target energy infimum over the
reachable scaffold class.  There exists \(\mu_{\mathrm{SQGI}}>0\) such
that every nonterminal selector step satisfies
\[
\Xi_T
\ge
\mu_{\mathrm{SQGI}}(E_T-E_\star).
\]
Assume additionally that
\[
\nu:=
\frac{c_{\mathrm{dec}}\mu_{\mathrm{SQGI}}}{C_K}
\in(0,1].
\]
\end{schurassumption}

\begin{schurlemma}[Perturbed linear recursion]
Let \(d_T\ge0\) and suppose
\[
d_{T+1}\le(1-\nu)d_T+\varepsilon_T
\]
for some \(\nu\in(0,1]\) and some nonnegative summable sequence
\(\{\varepsilon_T\}_{T\ge0}\).  Then
\[
d_T
\le
(1-\nu)^T d_0
+\sum_{s=0}^{T-1}(1-\nu)^{T-1-s}\varepsilon_s,
\]
and the convolution term tends to zero.  Hence \(d_T\to0\).
\end{schurlemma}

\begin{proof}
Iterating the recursion gives
\[
d_T
\le
(1-\nu)^T d_0
+\sum_{s=0}^{T-1}(1-\nu)^{T-1-s}\varepsilon_s.
\]
The first term tends to zero.  For the second term, fix
\(\epsilon>0\).  Since \(\sum_s\varepsilon_s<\infty\), choose \(T_0\)
such that \(\sum_{s=T_0}^{\infty}\varepsilon_s<\epsilon\).  Split the
convolution into
\[
\sum_{s=0}^{T-1}(1-\nu)^{T-1-s}\varepsilon_s
=
\sum_{s=0}^{T_0-1}(1-\nu)^{T-1-s}\varepsilon_s
+\sum_{s=T_0}^{T-1}(1-\nu)^{T-1-s}\varepsilon_s.
\]
The first sum is finite and tends to zero as \(T\to\infty\); the second
is bounded by \(\sum_{s=T_0}^{\infty}\varepsilon_s<\epsilon\).  Since
\(\epsilon\) was arbitrary, the convolution tends to zero.
\end{proof}

\begin{schurtheorem}[Conditional global convergence]
Assume the hypotheses of the actual-energy convergence theorem and the
Schur-QGI condition.  Then
\[
E_T\to E_\star.
\]
More precisely,
\[
E_T-E_\star
\le
(1-\nu)^T(E_0-E_\star)
+e_T,
\]
where
\[
e_T:=
\sum_{s=0}^{T-1}
(1-\nu)^{T-1-s}\varepsilon_s^{\mathrm{refit}},
\qquad
e_T\to0.
\]
In the exact-refit case
\(\varepsilon_T^{\mathrm{refit}}\equiv0\), this reduces to the geometric
bound
\[
E_T-E_\star
\le
(1-\nu)^T(E_0-E_\star).
\]
\end{schurtheorem}

\begin{proof}
By actual-energy descent and the selected-gain lower bound,
\[
\begin{aligned}
E_{T+1}
&\le
E_T
-\frac{\c_{\mathrm{dec}}}{C_K}\Xi_T
+\varepsilon_T^{\mathrm{refit}}.
\end{aligned}
\]
Using Schur-QGI,
\[
E_{T+1}
\le
E_T
-\frac{\c_{\mathrm{dec}}\mu_{\mathrm{SQGI}}}{C_K}
(E_T-E_\star)
+\varepsilon_T^{\mathrm{refit}}.
\]
With \(d_T:=E_T-E_\star\) and
\(\nu=c_{\mathrm{dec}}\mu_{\mathrm{SQGI}}/C_K\), this is
\[
d_{T+1}\le(1-\nu)d_T+\varepsilon_T^{\mathrm{refit}}.
\]
The perturbed linear recursion lemma gives the stated bound and
\(d_T\to0\).
\end{proof}

\begin{schurremark}[Interpretation]
The local convergence theorem does not require diagonal dominance of a
pool metric and does not require inversion of a raw, possibly singular
pool covariance matrix.  The role of diagonal dominance is replaced by
the computable reduced Schur residual
\[
F_{r,\mathrm{safe}}^{\mathrm{red}}\mathcal N_3(r;T)
=
\left[
F_{r,\mathrm{safe}}^{\mathrm{red}}
-
(q_r^{\mathrm{red}})^\top
(Q_r+\varepsilon_{\mathrm{nov}}I)^{-1}
q_r^{\mathrm{red}}
\right]_{+}.
\]
The price of claiming global convergence is the Schur-QGI condition:
the best surviving reduced-Schur gain must dominate the current
suboptimality \(E_T-E_\star\).  Without that condition, the theorem
still gives actual-energy convergence and vanishing Phase-3
Schur-stationarity residual \(\Xi_T\to0\), but not necessarily global
optimality.
\end{schurremark}

\fi
\end{document}
